% Gauss Werke Band 1: Disquisitiones Arithmeticae
% Sectio Quinta: De Formis Secundi Gradus
% Articles 165-200 (pages 141-190 of original)
\documentclass[12pt,a4paper]{article}
\usepackage[utf8]{inputenc}
\usepackage[T1]{fontenc}
\usepackage[latin,ngerman]{babel}
\usepackage{amsmath,amssymb,amsthm}
\usepackage{geometry}
\usepackage{hyperref}
\usepackage{enumitem}
\usepackage{mathrsfs}
\usepackage{longtable}
\usepackage{booktabs}
\usepackage{microtype}
\usepackage{yfonts}
\geometry{margin=2.5cm}

\theoremstyle{plain}
\newtheorem*{satz}{Satz}
\newtheorem*{theorem}{Theorem}

\pagestyle{headings}

\title{Werke Band 1 --- Disquisitiones Arithmeticae\\[0.5em]
\large Sectio Quinta. De Formis Secundi Gradus.\\
Articuli 165--200}
\author{Carolus Fridericus Gauss}
\date{}

\begin{document}
\maketitle

\selectlanguage{latin}

\section*{Generalia de repraesentationibus numerorum}

\textbf{165.}

Forma $9x^2 + 14xy - 4yy$ in formam $-12x'^2 - 18x'y' + 39y'^2$
transmutatur tum proprie, ponendo
\[
x = 4x' + 11y', \quad y = -2y',
\]
tum improprie, ponendo
\[
x = -74x' + 89y', \quad y = 15x' - 18y'.
\]

Hic igitur $a + a'$, $b + b'$, $a'' + b''$, $c + c'$ sunt $-70$, $100$, $-20$; autem $r'$, $s'$, $t'$, $u'$, $v'$ sunt $14$, $3$, $14$, $5$, $5$, $1$, $=1$.

$\mu = 0$, $\nu = -1$. Numeri autem $a$, $b$, $c$ inveniuntur $-237$, $-1170$, $48$, quorum divisor communis maximus $= 3 = r$; denique $\varepsilon = 3$.

Hinc transformatio $(S)$ fit haec, $x = bt - u$, $y = -at$. Per quam forma $(3,7,-4)$ transit in formam ancipitem $tt - 10tu + uu$.

\medskip

Si formae $F$, $F'$ sunt aequivalentes: forma $G$, sub $F$ contenta, etiam sub $F'$ contenta erit. Sed quoniam eandem formam etiam implicat, ipsi aequivalens erit, et proin etiam formae $F$. In hoc igitur casu theorema ita enunciabitur:

\begin{quote}
Si $F$, $F'$ tam proprie, quam improprie sunt aequivalentes: forma anceps utrique aequivalens inveniri poterit.
\end{quote}

Ceterum in hoc casu $\varepsilon = +1$, adeoque etiam $r$, ipsum $\varepsilon$ metiens, $=1$ erit.

\medskip

Haec de formarum transformatione in genere sufficiant: transimus itaque ad considerationem repraesentationum.

\bigskip

\begin{center}
\Large Generalia de repraesentationibus numerorum per formas, earumque nexu cum transformationibus.
\end{center}

\bigskip

\textbf{166.}

Si forma $F$ formam $F'$ implicat: quicunque numerus per $F'$ repraesentari potest, etiam per $F$ poterit.

Sint indeterminatae formarum $F$, $F'$ respective $x,y$; $x',y'$: ponamusque numerum $M$ per $F'$ repraesentari, faciendo $x'=m$, $y'=n$, formam $F$ vero in $F'$ transire per substitutionem
\[
x = \alpha x' + \beta y', \quad y = \gamma x' + \delta y'.
\]

Tum manifestum est, si ponatur
\[
x = \alpha m + \beta n, \quad y = \gamma m + \delta n,
\]
$F$ transire in $M$.

Si $M$ pluribus modis per formam $F'$ repraesentari potest, e.g.~etiam faciendo $x'=m'$, $y'=n'$: plures repraesentationes ipsius $M$ per $F$ inde sequentur. Si enim esset tum
\[
\alpha m + \beta n = \alpha m' + \beta n' \quad \text{tum} \quad \gamma m + \delta n = \gamma m' + \delta n'
\]
foret aut $\alpha\delta - \beta\gamma = 0$ adeoque etiam determinans formae $F = 0$ contra hyp.,

aut $m = m'$, $n = n'$. Hinc sequitur, $M$ ad minimum totidem modis diversis per $F$ repraesentari posse, quot per $F'$.

Si igitur tum $F$ ipsam $F'$, tum $F'$ ipsam $F$ implicat i.e.~si $F$, $F'$ sunt aequivalentes: numerusque $M$ per alterutram repraesentari potest: etiam per alteram repraesentari poterit, et quidem totidem modis diversis per alteram, quot per alteram.

Denique observamus, in hocce casu divisorem communem maximum numerorum $m,n$ aequalem esse divisori comm.~max.~numerorum $\alpha m + \beta n$, $\gamma m + \delta n$. Sit ille $= \mu$, numerique $\mu$, $\nu$ ita accepti, ut fiat $\mu m + \nu n = \mu$. Tum erit
\[
(\delta\mu - \gamma\nu)(\alpha m + \beta n) - (\beta\mu - \alpha\nu)(\gamma m + \delta n) = (\alpha\delta - \beta\gamma)(\mu m + \nu n) = \pm \mu.
\]
Hinc div.~comm.~max.~numerorum $\alpha m + \beta n$, $\gamma m + \delta n$ metietur ipsum $\mu$, $\mu$ vero etiam illum metietur, quia manifeste ipsos $\alpha m + \beta n$, $\gamma m + \delta n$ metitur. Quare necessario ille erit $= \mu$. Quando igitur $m$, $n$ inter se primi sunt, etiam $\alpha m + \beta n$, $\gamma m + \delta n$ inter se primi erunt.

\bigskip

\textbf{167.}

\textbf{Theorem.} Si formae
\begin{align*}
axx + 2bxy + cyy \tag{$F$}\\
a'x'x' + 2b'x'y' + c'y'y' \tag{$F'$}
\end{align*}
sunt aequivalentes ipsarum determinans $= D$, posteriorque in priorem transit ponendo
\[
x = \alpha x' + \beta y', \quad y = \gamma x' + \delta y';
\]
porro numerus $M$ per $F$ repraesentatur faciendo $x = m$, $y = n$, adeoque per $F'$ faciendo
\[
x' = \alpha m + \beta n = m', \quad y' = \gamma m + \delta n = n'
\]
et quidem ita, ut $m$ ad $n$ eoque ipso etiam $m'$ ad $n'$ sit primus: ambae repraesentationes aut ad eundem valorem expressionis $\sqrt{D} \pmod{M}$ pertinebunt, aut ad oppositos, prout transformatio formae $F'$ in $F$ propria est vel impropria.

\textit{Dem.} Determinentur numeri $\mu$, $\nu$ ita, ut fiat $\mu m + \nu n = 1$, ponaturque
\[
\mu(mb+nc) - \nu(ma+nb) = N,
\]
qui erunt integri propter $\alpha\delta - \beta\gamma = \pm 1$. Tum erit $\mu m + \nu n = 1$. (Cf.~art.~prae.~fin.)

Porro sit
\[
\mu(b'm' + c'n') - \nu(a'm' + b'n') = N';
\]
eruntque $N$, $N'$ valores expr.~$\sqrt{D} \pmod{M}$, ad quos repraesentatio prima et secunda pertinent.

Si in $N$ pro $a'$, $b'$, $c'$ valores ipsorum substituuntur,
\begin{align*}
\text{pro } a, \quad & \alpha\alpha a' + 2\alpha\gamma b' + \gamma\gamma c',\\
\text{pro } b, \quad & \alpha\beta a' + b'(\alpha\delta + \beta\gamma) + c'\gamma\delta,\\
\text{pro } c, \quad & \beta\beta a' + 2\beta\delta b' + \delta\delta c',
\end{align*}
invenietur evolutione facta $N = N'(\alpha\delta - \beta\gamma)$.

Quare erit aut $N = N'$, aut $N = -N'$, prout $\alpha\delta - \beta\gamma = +1$ aut $=-1$, i.e.~repraesentationes pertinebunt ad eundem valorem expr.~$\sqrt{D} \pmod{M}$ vel ad oppositos, prout transformatio formae $F'$ in $F$ est propria vel impropria. \hfill Q.E.D.

Si itaque plures repraesentationes numeri $M$ per formam $(a,b,c)$ ope valorum inter se primorum indeterminatarum $x,y$, habentur ad valores diversos expr.~$\sqrt{D} \pmod{M}$ pertinentes: repraesentationes respondentes per formam $(a',b',c')$ ad eosdem resp.~valores pertinebunt, et si nulla repraesentatio numeri $M$ per formam aliquam ad valorem quendam determinatum pertinens datur, nulla quoque dabitur ad hunc valorem pertinens per formam illi aequivalentem.

\bigskip

\textbf{168.}

\textbf{Theorem.} Si numerus $M$ per formam $axx + 2bxy + cyy$ repraesentatur, tribuendo ipsis $x,y$ valores inter se primos $m,n$, valorque expressionis $\sqrt{D} \pmod{M}$, ad quem haec repraesentatio pertinet, est $N$: formae $(a,b,c)$, $(M,N,\frac{NN-D}{M})$ proprie aequivalentes erunt.

\textit{Demonstr.} Ex art.~155 patet, numeros integros $\mu$, $\nu$ inveniri posse ita ut sit
\[
\mu m + \nu n = 1, \quad \mu(mb+nc) - \nu(ma+nb) = N.
\]

Quo facto forma $(a,b,c)$ per substitutionem $x = mx' - \nu y'$, $y = nx' + \mu y'$, quae manifesto est propria, transit in formam, cuius determinans $= D(m\mu + n\nu)^2 = D$, sive in formam aequivalentem: quae forma si ponitur $= (M', N', \frac{N'N'-D}{M'})$, erit
\[
M' = amm + 2bmn + cnn = M, \quad N' = -m\mu a + (m\nu - n\mu)b + n\nu c = N.
\]
Quare forma, in quam $(a,b,c)$ per transformationem illam mutatur, erit $(M,N,\frac{NN-D}{M})$. \hfill Q.E.D.

Ceterum ex aequationibus
\[
m\mu + n\nu = 1, \quad \mu(mb+nc) - \nu(ma+nb) = N
\]
deducitur
\[
\mu = \frac{nN + ma + nb}{M}, \quad \nu = \frac{mb + nc - mN}{M},
\]
qui numeri itaque erunt integri.

Porro observandum, hanc propositionem locum non habere, si $M = 0$; tum enim terminus $\frac{NN-D}{M}$ fit indeterminatus\footnote{In hoc enim casu, si ad ipsum phrasin extendere volumus, haec: $N$ esse valorem expr.~$\sqrt{D} \pmod{M}$, sive $NN \equiv D \pmod{M}$ significabit $NN-D$ esse multiplum ipsius $M$, adeoque $= 0$.}.

\bigskip

\textbf{169.}

Si plures repraesentationes numeri $M$ per $(a,b,c)$ habentur, ad eundem valorem expr.~$\sqrt{D} \pmod{M}$, $N$, pertinentes (ubi valores ipsorum $x,y$ semper inter se primos supponimus): plures etiam transformationes propriae formae $(a,b,c)\ldots(F)$ in $(M,N,\frac{NN-D}{M})\ldots(G)$ inde deducentur. Scilicet si etiam per hos valores $x = m$, $y = n$ talis repraesentatio provenit, $(F)$ etiam per substitutionem
\[
x = mx' + \frac{m'N - m'b - n'c}{M}y', \quad y = nx' + \frac{n'N + m'a + n'b}{M}y'
\]
in $(G)$ transit. Vice versa, ex quavis transformatione propria formae $(F)$ in $(G)$ sequetur repraesentatio numeri $M$ per formam $(F)$, ad valorem $N$ pertinens. Scilicet si $(F)$ transit in $(G)$ positis $x = mx' + \mu y'$, $y = nx' + \nu y'$, $M$ repraesentatur per $(F)$ ponendo $x = m$, $y = n$, et quoniam hic $m\mu + n\nu = 1$ valor expr.~$\sqrt{D} \pmod{M}$, ad quem repraesentatio pertinet, erit $\mu(bm+cn) - \nu(am+bn)$ i.e.~$N$. Ex pluribus vero transformationibus propriis diversis sequentur totidem repraesentationes diversae ad $N$ pertinentes\footnote{Si ex duabus transformationibus propriis diversis eadem repraesentatio defluere supponitur, illae ita se habere debebunt 1) $x = mx' - \nu y'$, $y = nx' + \mu y'$; 2) $x = mx' - \nu' y'$, $y = nx' + \mu' y'$. Sed ex duabus aequationibus $m\mu + n\nu = m\mu' + n\nu'$, $\mu(mb+nc) - \nu(ma+nb) = \mu'(mb+nc) - \nu'(ma+nb)$ facile deducitur, esse aut $M = 0$ aut $m = m'$, $n = n'$, $\mu = \mu'$, $\nu = \nu'$. Casum $M = 0$ iam exclusimus.}. Hinc facile colligitur, si omnes transformationes propriae formae $(F)$ in $(G)$ habeantur, ex his omnes repraesentationes ipsius $M$ per $(F)$ ad valorem $N$ pertinentes sequi. Unde quaestio de repraesentationibus numeri dati per formam datam (in quibus indeterminatae valores inter se primos nanciscuntur) investigandis, reducta est ad quaestionem de inveniendis omnibus transformationibus propriis formae illius in datam aequivalentem.

Applicando iam ad haec ea, quae in art.~162 docuimus, facile concluditur:

Si repraesentatio aliqua numeri $M$ per formam $(F)$ ad valorem $N$ pertinens sit haec: $x = \alpha$, $y = \gamma$: formulam generalem omnes repraesentationes eiusdem numeri per formam $(F)$, ad valorem $N$ pertinentes, comprehendentem fore hanc:
\[
x = \alpha t - (\alpha h + \gamma c)u, \quad y = \gamma t + (\alpha a + \gamma h)u,
\]
umi $m$ divisor communis maximus numerorum $a$, $2b$, $c$; et $t$, $u$ omnes numeri indefinite, aequationi $tt - Duu = mm$ satisfacientes.

Si forma $(a,b,c)$ ancipiti alicui aequivalens, adeoque formae $(M,N,\frac{NN-D}{M})$ tam proprie, quam improprie, sive tam formae $(M,N,\frac{NN-D}{M})$, quam huic $(M,-N,\frac{NN-D}{M})$ proprie: repraesentationes numeri $M$ habebuntur per formam $(F)$, tam ad valorem $N$, quam ad valorem $-N$ pertinentes. Et vice versa si plures repraesentationes numeri $M$ per eandem formam $(F)$ ad valores oppositos expr.~$\sqrt{D} \pmod{M}$, $N$, $-N$, pertinentes habentur: forma $(F)$ formae $(G)$ tam proprie quam improprie aequivalens erit, formaque anceps assignari poterit, cui $(F)$ aequivaleat.

Haec generalia de repraesentationibus hic sufficiant: de repraesentationibus, in quibus indeterminatae valores inter se non primos habent, infra dicemus. Respectu aliarum proprietatum, formae, quarum determinans est negativus, prorsus alio modo quam formae determinantis positivi sunt tractandae: quare iam utrasque seorsim considerabimus. Ab illis tamquam facilioribus initium facimus.

\bigskip

\begin{center}
\Large De formis determinantis negativi.
\end{center}

\bigskip

\textbf{171.}

\textbf{Problema.} Proposita forma quacunque, $(a,b,a')$, cuius determinans negativus, $= -D$, designante $D$ numerum positivum, invenire formam huic proprie aequivalentem, $(A,B,C)$, in qua $A$ nec maior quam $\sqrt{\frac{4}{3}D}$, $C$, nec minor quam $2B$.

\textit{Solutio.} Supponimus in forma proposita non omnes tres conditiones simul locum habere; alioquin enim aliam formam quaerere opus non esset. Sit $b'$ residuum abs.~min.~numeri $-b$ secundum modulum $a'$, atque $a'' = \frac{bb+D}{a'}$; erit integer, quia $b'b' \equiv bb$, $b'b'+D \equiv bb+D = aa' \pmod{a}$. Iam si $a'' < a'$, fiat denuo $b''$ resid.~abs.~min.~ipsius $-b'$ secundum mod.~$a''$, atque $a''' = \frac{b'b'+D}{a''}$. Si hic iterum $a''' < a''$, sit rursus $b'''$ res.~abs.~min.~ipsius $-b''$ secundum mod.~$a''$ atque $a''' = \frac{b''b''+D}{a''}$. Haec operatio continuetur, donec in progressione $a'$, $a''$, $a'''$, $a^{IV}$ etc.~ad terminum $a^m$ perveniatur qui praecedente suo $a^{m-1}$ non sit minor, quod tandem evenire debet, quia alias progressio infinita numerorum integrorum continuo decrescentium haberetur. Tum forma $(a^m, b^m, a^{m+1})$ omnibus conditionibus satisfaciet.

\textit{Dem.} I. In progressione formarum $(a,b,a')$, $(a',b',a'')$, $(a'',b'',a''')$ etc.~quaevis praecedenti est contigua, quare ultima primae proprie aequivalens erit (artt.~159,~160).

II. Quum $b^m$ sit residuum absolute minimum ipsius $-b^{m-1}$ secundum mod.~$a^m$, maior quam $\frac{1}{2}a^m$ non erit (art.~4).

III. Quia $a^m a^{m+1} = D + b^m b^m$, atque $a^{m+1}$ non $< a^m$, $a^m a^{m+1}$ non erit $> D + b^m b^m$, et quum $b^m$ non $> \frac{1}{2}a^m$, $a^m a^{m+1}$ non erit $> D + \frac{1}{4}a^m a^m$ et $a^m a^{m+1}$ non $> D$, tandemque $a^m$ non $> \sqrt{\frac{4}{3}D}$.

\textit{Exempl.} Proposita sit forma $(304, 217, 155)$, cuius determinans $= -31$.

Hic invenitur progressio formarum
\[
(304, 217, 155), \quad (155, -62, 25), \quad (25, 12, 7), \quad (7, 2, 5), \quad (5, -2, 7).
\]
Ultima est quaesita. Eodem modo formae $(121, 49, 20)$, cuius determinans $= -19$, aequivalentes inveniuntur: $(20, 9, 5)$, $(5, -1, 4)$, $(4, 1, 5)$: quare $(4, 1, 5)$ erit forma quaesita.

Tales formas $(A, B, C)$, quarum determinans est negativus et in quibus $A$ nec maior quam $\sqrt{\frac{4}{3}D}$, $C$ nec minor quam $2B$, \textbf{formas reductas} vocabimus. Quare cuivis formae determinantis negativi, forma reducta proprio aequivalens inveniri poterit.

\bigskip

\textbf{172.}

\textbf{Problema.} Invenire conditiones, sub quibus duae formae reductae non identicae, eiusdem determinantis $-D$, $(a,b,c)$, $(a',b',c')$ proprie aequivalentes esse possint.

\textit{Solutio.} Supponamus, id quod licet, $a$ esse non $> c$, formamque $axx + 2bxy + cyy$ transire in $a'x'x' + 2b'x'y' + c'y'y'$ per substitutionem propriam $x = \alpha x + \beta y$, $y = \gamma x + \delta y$. Tum habebuntur aequationes
\begin{align*}
a\alpha\alpha + 2b\alpha\gamma + c\gamma\gamma &= a' \tag{1}\\
a\alpha\beta + b(\alpha\delta + \beta\gamma) + c\gamma\delta &= b' \tag{2}\\
\alpha\delta - \beta\gamma &= 1 \tag{3}
\end{align*}

Ex (1) sequitur $a'a = (a\alpha + b\gamma)^2 + D\gamma\gamma$; quare $a'a$ erit positivus; et quum $ac = D + bb$, $a'c' = D + b'b'$, etiam $ac$, $a'c'$ positivi erunt: quare $a$, $a'$, $c$, $c'$ omnes eadem signa habebunt. Sed tum $a$ tum $a'$ non $> \sqrt{\frac{4}{3}D}$, adeoque $a'a'$ non $> \frac{4}{3}D$; quare multo minus $D\gamma\gamma (= a'a - (a\alpha + b\gamma)^2)$ maior quam $\frac{4}{3}D$ esse poterit. Hinc $\gamma$ erit aut $=0$, aut $= \pm 1$.

\textbf{I.} Si $\gamma = 0$, ex (3) sequitur, esse aut $\alpha = 1$, $\delta = 1$, aut $\alpha = -1$, $\delta = -1$.

In utroque casu fit ex (1) $a = a'$, et ex (2) $b' - b = \pm a\alpha\beta$. Sed $b$ non $> \frac{1}{2}a$, et $b'$ non $> \frac{1}{2}a$ proin etiam non $> \frac{1}{2}a$. Quare aequatio $b' - b = \pm a\alpha\beta$ consistere nequit, nisi fuerit

aut $b = b'$, unde sequeretur $c = \frac{D+bb}{a} = \frac{D+b'b'}{a'} = c'$, quare formae $(a,b,c)$, $(a',b',c')$ identicae essent contra hyp.

aut $b = -b' \pm \frac{1}{2}a$. In hoc etiam casu erit $c = c'$ formaque $(a',b',c')$ erit $(a,-b,c)$ i.e.~formae $(a,b,c)$ opposita. Simul patet, formas has esse ancipites propter $2b = \pm a$.

\textbf{II.} Si $\gamma = \pm 1$, fit ex (1) $a\alpha\alpha + c = a' + 2b\alpha$. Sed $c$ non minor quam $a$, adeoque non minor quam $a'$; hinc $a\alpha\alpha + c - a'$ sive $2b\alpha$ certo non minor quam $a\alpha\alpha$. Quare quum $2b$ non sit maior quam $a$, erit $\alpha\alpha$ non minor quam $\alpha$; unde necessario aut $\alpha = 0$, aut $\alpha = \pm 1$.

1) Si $\alpha = 0$, fit ex (1) $a' = c$, et quoniam $a$ neque maior quam $c$, neque minor quam $a'$, erit necessario $a' = a = c$. Porro ex (3) fit $\beta\gamma = -1$, unde ex (2) $b + b' = \pm \frac{1}{2}c = \pm \frac{1}{2}a$. Hinc simili modo ut in (I) sequitur esse

aut $b = b'$, in quo casu formae $(a,b,c)$, $(a',b',c')$ forent identicae, contra hyp.

aut $b = -b'$, in quo casu formae $(a',b',c')$ erunt oppositae $(a,b,c)$.

2) Si $\alpha = \pm 1$, ex (1) sequitur $+2b = \alpha(a' + c - a)$. Quare quum neque $a$ neque $c < a'$, erit $2b$ non $< a$ et non $< c$. Sed $2b$ etiam non $> a$ neque $> c$, unde necessario $\pm 2b = a = c$, et hinc ex aequ.~$2b = a + c - a'$, etiam $= a'$. Fit igitur ex (2)
\[
b' = a(\alpha\beta + \gamma\delta) + b(\alpha\delta + \beta\gamma)
\]
sive, propter $\alpha\delta - \beta\gamma = 1$,
\[
b' - b = a(\alpha\beta + \gamma\delta \pm \beta\gamma).
\]
quare necessario, ut ante

aut $b = b'$, unde formae $(a,b,c)$, $(a',b',c')$ identicae, contra hyp.

aut $b = -b'$, adeoque formae illae oppositae. Simul in hoc casu propter $a = \pm 2b$, formae erunt ancipites.

\medskip

Ex his omnibus colligitur, formas $(a,b,c)$, $(a',b',c')$ proprie aequivalentes esse non posse, nisi fuerint oppositae, simulque aut ancipites, aut $a = c = a' = c'$.

In hisce casibus formas $(a,b,c)$, $(a',b',c')$ proprie aequivalere, vel a priori facile praevideri potuit; si enim formae sunt oppositae, improprie, et si insuper ancipites, etiam proprie aequivalentes esse debent; si vero $a = c$, forma $\frac{D+(a-b)^2}{2a}$ formae $(a,b,c)$ contigua et proin aequivalens erit; sed propter $D + bb = ac = aa$ fit $\frac{D+(a-b)^2}{2a} = \frac{2aa - 2ab}{2a} = 2a - 2b$, forma vero $(2a-2b, a-b, a)$ est anceps; quare $(a,b,c)$ oppositae suae etiam proprie aequivalebit.

Aeque facile iam diiudicari potest, quando duae formae reductae $(a,b,c)$, $(a',b',c')$ non oppositae improprie aequivalentes esse possunt. Erunt enim impr. aequivalentes, si $(a,b,c)$, $(a',-b',c')$, quae non identicae erunt, proprie sunt aequivalentes, et contra. Hinc patet, conditionem, sub qua illae improprie sint aequivalentes, esse, ut sint identicae, insuperque aut ancipites aut $a = c$.

Formae vero reductae, quae neque identicae sunt neque oppositae, neque proprie neque improprie aequivalentes esse possunt.

\bigskip

\textbf{173.}

\textbf{Problema.} Propositis duabus formis eiusdem determinantis negativi, $F$ et $F'$, investigare, utrum sint aequivalentes.

\textit{Solutio.} Quaerantur duae formae reductae $f$, $f'$ formis $F$, $F'$ resp.~proprie aequivalentes: si formae $f$, $f'$ sunt proprie, vel improprie vel utroque modo aequivalentes, etiam $F$, $F'$ erunt; si vero $f$, $f'$ nullo modo aequivalentes sunt, etiam $F$, $F'$ non erunt.

Ex art.~prae.~dari possunt quatuor casus:

1) Si $f$, $f'$ neque identicae neque oppositae, $F$, $F'$ nullo modo aequivalentes erunt.

2) Si $f$, $f'$ sunt primo vel identicae vel oppositae, et secundo vel ancipites, vel terminos suos extremos aequales habent: $F$, $F'$ tum proprie, tum improprie aequivalentes erunt.

3) Si $f$, $f'$ sunt identicae, neque vero ancipites neque terminos extremos aequales habent: $F$, $F'$ proprie tantum aequivalebunt.

4) Si $f$, $f'$ sunt oppositae, neque vero ancipites, neque terminos extremos aequales habent: $F$, $F'$ improprie tantum aequivalentes erunt.

\textit{Ex.} Formis $(41,35,30)$, $(7,18,47)$, quarum determinans $= -5$, reductae $(1,0,5)$, $(2,1,3)$ non aequivalentes inveniuntur, quare illae nullo modo aequivalentes erunt. Formis vero $(23,38,63)$, $(15,20,27)$ aequivalet eadem reducta $(2,1,3)$, quae quum simul sit anceps, formae $(23,38,63)$, $(15,20,27)$ tum proprie tum improprie aequivalebunt. Formis $(37,53,78)$, $(53,73,102)$, aequivalent reductae $(9,2,9)$, $(9,-2,9)$, quae quum sint oppositae, ipsarumque termini extremi aequales: formae propositae tam proprie quam improprie erunt aequivalentes.

\bigskip

\textbf{174.}

Multitudo omnium formarum reductarum, determinantem datum $-D$ habentium, semper est finita, et, respectu numeri $D$, satis modica: formae hae ipsae vero duplici modo inveniri possunt. Designemus formas reductas determinantis $-D$ indefinite per $(a,b,c)$, ubi itaque omnes valores ipsorum $a$, $b$, $c$ determinari debent.

\textbf{Methodus prima.} Accipiantur pro $a$ omnes numeri, tum positivi tum negativi non maiores quam $\sqrt{\frac{4}{3}D}$, quorum residuum quadraticum $= -D$, et pro singulis $a$, fiat $b$ successive aequalis omnibus valoribus expr.~$\sqrt{-D} \pmod{a}$, maioribus quam $\frac{1}{2}a$, tum positive tum negative acceptis; $c$ vero pro singulis valoribus determinatis ipsorum $a$, $b$, ponatur $= \frac{bb+D}{a}$. Si quae formae hoc modo oriuntur, in quibus $c < a$, hae erunt reiiciendae, reliquae autem manifesto erunt reductae.

\textbf{Methodus secunda.} Accipiantur pro $b$ omnes numeri tum positivi tum negativi, non maiores quam $\frac{1}{2}\sqrt{\frac{4}{3}D}$, sive $\frac{1}{3}\sqrt{3D}$; pro singulis $b$ resolvatur $bb+D$ omnibus quibus fieri potest modis in binos factores (etiam signorum diversitatis ratione habita) ambos ipso $2b$ non minores ponaturque alter factor, et quidem, quando factores sunt inaequales, minor $= a$, alter $= c$. Quum $a$ non erit $> \frac{2}{3}\sqrt{3D}$, omnes formae, quae hoc modo prodeunt, manifesto erunt reductae. Denique patet, nullam formam reductam dari posse, quae non per utramque methodum inveniatur.

\textit{Ex.} Sit $D = 85$. Hic limes valorum ipsius $a$ est $\sqrt{\frac{340}{3}}$, qui iacet inter $10$ et $11$. Numeri vero inter $1$ et $10$ (incl.), quorum residuum $= -85$, sunt $1$, $2$, $5$, $10$. Unde habentur formae duodecim:
\[
\begin{aligned}
&(1,0,85), &&(2,1,43), &&(2,-1,43), &&(5,0,17), &&(10,5,11), &&(10,-5,11);\\
&(-1,0,-85), &&(-2,1,-43), &&(-2,-1,-43), &&(-5,0,-17), &&(-10,5,-11), &&(-10,-5,-11).
\end{aligned}
\]
Per methodum alteram limes valorum ipsius $b$ habetur $\frac{1}{3}\sqrt{255}$, qui situs est inter $5$ et $6$. Pro $b = 0$, prodeunt formae
\[
(1,0,85), \quad (-1,0,-85), \quad (5,0,17), \quad (-5,0,-17),
\]
pro $b = \pm 1$ hae $(2,1,43)$, $(-2,-1,-43)$.

Pro $b = \pm 2$ nullae habentur, quia $89$ in duos factores, qui ambo non $< 4$, resolvi nequit. Idem valet de $\pm 3$, $\pm 4$. Tandem pro $b = \pm 5$, proveniunt
\[
(10,5,11), \quad (-10,-5,-11).
\]

\bigskip

\textbf{175.}

Si ex omnibus formis reductis determinantis dati, formarum binarum, quae, licet non identicae, tamen proprie sunt aequivalentes, alterutra reiicitur; formae remanentes hac insigni proprietate erunt praeditae, ut, quaevis forma eiusdem determinantis alicui ex ipsis proprie sit aequivalens, et quidem unicae tantum (alias enim inter ipsas aliquae proprie aequivalentes forent). Unde patet, omnes formas eiusdem determinantis in totidem classes distribui posse, quot formae remanserint, referendo scilicet formas eidem reductae proprie aequivalentes in eandem classem. Ita pro $D = 85$, remanent formae
\[
\begin{aligned}
&(1,0,85), &&(2,1,43), &&(5,0,17), &&(10,5,11),\\
&(-1,0,-85), &&(-2,1,-43), &&(-5,0,-17), &&(-10,5,-11)
\end{aligned}
\]
quare omnes formae determinantis $-85$ in octo classes distribui poterunt, prout formae primae, aut secundae etc.~proprie aequivalent. Perspicuum vero est, formas in eadem classe locatas proprie aequivalentes fore, formas ex diversis classibus proprie aequivalentes esse non posse. Sed hoc argumentum de classificatione formarum infra multo fusius exsequemur.

Hic unicam observationem adiicimus. Iam supra ostendimus, si determinans formae $(a,b,c)$ fuerit negativus $= -D$, $a$ et $c$ eadem signa habere (quia scilicet $ac = bb + D$ adeoque positivus); eadem ratione facile perspicitur si formae $(a,b,c)$, $(a',b',c')$ sint aequivalentes, omnes $a$, $c$, $a'$, $c'$ eadem signa habituros. Si enim prior in posteriorem per Substitut. $x = \alpha x' + \beta y'$, $y = \gamma x' + \delta y'$ transit: erit $a\alpha\alpha + 2b\alpha\gamma + c\gamma\gamma = a'$, hinc $a'a = (a\alpha + b\gamma)^2 + D\gamma\gamma$ adeoque certo non negativus; quoniam vero neque $a$ neque $a' = 0$ esse potest, erit $aa'$ positivus et proin signa ipsorum $a$, $a'$ eadem. Hinc manifestum est, formas, quarum termini exteri sint positivi, ab iis, quorum termini exteri sint negativi, prorsus esse separatas, sufficitque ex formis reductis eas tantum considerare, quae terminos suos exteros positivos habent, nam reliquae totidem sunt multitudine et ex illis oriuntur, tribuendo terminis exteris signa opposita; idemque valet de formis ex reductis reiiciendis et remanentibus.

\bigskip

\textbf{176.}

Ecce itaque pro determinantibus quibusdam negativis tabulam formarum, secundum quas omnes reliquae eiusdem determinantis in classes distingui possunt; apponimus autem, ad annotat.~art.~prae., semissem tantum, scilicet, eas quarum termini exteri positivi.

\begin{center}
\begin{tabular}{cl}
$D$ & \quad Formae reductae \\ \hline
\rule{0pt}{2.5ex}%
$1$ & $(1,0,1)$ \\
$2$ & $(1,0,2)$ \\
$3$ & $(1,0,3)$, $(2,1,2)$ \\
$4$ & $(1,0,4)$, $(2,0,2)$ \\
$5$ & $(1,0,5)$, $(2,1,3)$ \\
$6$ & $(1,0,6)$, $(2,0,3)$ \\
$7$ & $(1,0,7)$, $(2,1,4)$ \\
$8$ & $(1,0,8)$, $(2,0,4)$, $(3,1,3)$ \\
$9$ & $(1,0,9)$, $(2,1,5)$, $(3,0,3)$ \\
$10$ & $(1,0,10)$, $(2,0,5)$ \\
$11$ & $(1,0,11)$, $(2,1,6)$, $(3,1,4)$, $(3,-1,4)$ \\
$12$ & $(1,0,12)$, $(2,0,6)$, $(3,0,4)$, $(4,2,4)$ \\
\end{tabular}
\end{center}

Superfluum foret hanc tabulam hic ulterius continuare, quippe quam infra multo aptius disponere docebimus.

Patet itaque, quamvis formam determinantis $-1$, formae $xx + yy$ proprie aequivalere, si ipsius termini exteri sint positivi, vel huic $-xx - yy$, si sint negativi; quamvis formam determinantis $-2$, cuius termini exteri positivi, formae $xx + 2yy$ etc.; quamvis formam determinantis $-11$, cuius termini exteri positivi, alicui ex his $xx + 11yy$, $2xx + 2xy + 6yy$, $3xx + 2xy + 4yy$, $3xx - 2xy + 4yy$ etc.

\bigskip

\textbf{177.}

\textbf{Problema.} Habetur series formarum, quarum quaevis praecedenti a parte posteriori contigua: desideratur transformatio aliqua propria primae in formam quamcunque seriei.

\textit{Solutio.} Sint formae $(a,b,a') = F$; $(a',b',a'') = F'$; $(a'',b'',a''') = F''$; $(a''',b''',a^{IV}) = F'''$ etc. Designentur $\frac{b+b'}{a'}$, $\frac{b'+b''}{a''}$, $\frac{b''+b'''}{a'''}$ etc.~respective per $h'$, $h''$, $h'''$ etc. Sint indeterminatae formarum $F$, $F'$, $F''$ etc.~$x,y$; $x',y'$; $x'',y''$ etc.

Ponatur $F$ transmutari
\begin{align*}
\text{in } F' &\text{ positis } x = -y', \quad y = x' + h'y';\\
\text{in } F'' &\text{ positis } x = -y'', \quad y = x'' + h''y'';\\
\text{in } F''' &\text{ positis } x = -y''', \quad y = x''' + h'''y''' \text{ etc. (art. 160)}
\end{align*}
facile eruetur sequens algorithmus (art.~159):
\[
\begin{aligned}
\alpha' &= 0, & \beta' &= -1, & \gamma' &= 1, & \delta' &= h';\\
\alpha'' &= \gamma', & \beta'' &= h''\gamma' - \beta', & \gamma'' &= \alpha', & \delta'' &= h''\alpha' - \delta';\\
\alpha''' &= \gamma'', & \beta''' &= h'''\gamma'' - \beta'', & \gamma''' &= \alpha'', & \delta''' &= h'''\alpha'' - \delta'' \text{ etc.}
\end{aligned}
\]
sive
\[
\begin{aligned}
\alpha' &= 0, & \beta' &= -1;\\
\alpha'' &= 1, & \beta'' &= h'';\\
\alpha''' &= h''', & \beta''' &= h'''h'' + 1;\\
\alpha^{IV} &= h'''h'' + 1, & \beta^{IV} &= h^{IV}h'''h'' + h^{IV} + h'' \text{ etc.}
\end{aligned}
\]
haec ad formas determinantis negativi non est restricta, sed ad omnes casus patet, si modo nullus numerorum $a'$, $a''$, $a'''$ etc.~$= 0$.

\bigskip

\textbf{178.}

\textbf{Problema.} Propositis duabus formis $F$, $f$, eiusdem determinantis negativi, proprie aequivalentibus: invenire transformationem aliquam propriam alterius in alteram.

\textit{Sol.} Supponamus formam $F$ esse $(A, B, A')$ et per methodum art.~171 inventam esse progressionem formarum $(A, B, A')$, $(A', B', A'')$, $(A'', B'', A''')$ etc.~usque ad $(A^m, B^m, A^{m+1})$, quae sit reducta: similiterque $f$ esse $(a,b,a')$ et per eandem methodum inventam seriem $(a',b',a'')$, $(a'',b'',a''')$ usque ad $(a^n, b^n, a^{n+1})$, quae sit reducta. Tum duo casus locum habere possunt.

\textbf{I.} Si formae $(A^m, B^m, A^{m+1})$, $(a^n, b^n, a^{n+1})$ sunt aut identicae, aut oppositae simulque ancipites. Tum formae $(A^{m-1}, B^{m-1}, A^m)$, $(a^{n-1}, -b^{n-1}, a^n)$ erunt contiguae (designante $A^{m-1}$ terminum progressionis $A$, $A'$, $A''$\ldots$A^m$ penultimum, similiaque $a^{n-1}$, $b^{n-1}$, $a^n$). Nam $A^m \equiv B^{m-1} \equiv B^m \pmod{A^{m-1}}$, $b^{n-1} \equiv -b^n \pmod{a^n}$ sive $A^m$; unde $B^{m-1} - b^{n-1} \equiv B^m - B^m = 0$, si formae $(A^m, B^m, A^{m+1})$, $(a^n, b^n, a^{n+1})$ sunt identicae, et $\equiv 2B^m$ adeoque $= 0$, si sunt oppositae et ancipites. Quare in progressione formarum
\[
\begin{aligned}
&(A, B, A'), \quad (A', B', A''), \quad \ldots \quad (A^{m-1}, B^{m-1}, A^m),\\
&(a^{n-1}, -b^{n-1}, a^n), \quad (a^{n-2}, -b^{n-2}, a^{n-1}), \quad \ldots \quad (a, b, a'), \quad (a, b, a')
\end{aligned}
\]
quaevis forma praecedenti contigua erit, adeoque per art.~prae.~transformatio propria primae $F$ in ultimam $f$ inveniri poterit.

\textbf{II.} Si formae $(A^m, B^m, A^{m+1})$, $(a^n, b^n, a^{n+1})$ non identicae, sed oppositae simulque $A^m = A^{m+1} = a^n = a^{n+1}$. Tum progressio formarum
\[
\begin{aligned}
&(A, B, A'), \quad (A', B', A''), \quad \ldots \quad (A^{m-1}, B^{m-1}, A^m),\\
&(a^{n-1}, -b^{n-1}, a^n), \quad (a^{n-2}, -b^{n-2}, a^{n-1}), \quad \ldots \quad (a, b, a'), \quad (a, b, a')
\end{aligned}
\]
eadem proprietate erit praedita. Nam $A^{m+1} = a^n$, et $B^{m-1} - b^{n-1} = -(B^m + B^m)$ per $a^n$ divisibilis. Unde per art.~prae.~invenietur transformatio propria formae primae $F$ in ultimam $f$.

\textit{Ex.} Ita pro formis $(23, 38, 63)$, $(15, 20, 27)$ habetur progressio
\[
(23,38,63), \; (63,25,10), \; (10,5,3), \; (3,1,2), \; (2,-7,27), \; (27,-20,15), \; (15,20,27)
\]
quare
\[
h' = 1, \quad h'' = 3, \quad h''' = 2, \quad h^{IV} = -3, \quad h^V = -1, \quad h^{VI} = 3.
\]
Hinc deducitur transformatio formae $23xx + 76xy + 63yy$ in $15tt + 40tu + 27uu$ haec: $x = -13t - 18u$, $y = 8t + 11u$.

Ex solutione hac nullo negotio sequitur solutio problematis: Si formae $F$, $f$ improprie sunt aequivalentes, invenire transformationem impropriam formae $F$ in $f$. Sit enim $f = att + 2btu + duu$ eritque forma opposita $app - 2bpq + dqq$ formae $F$ proprio aequivalens. Quaeratur transformatio propria formae $F$ in illam, $X = \alpha p + \beta q$, $Y = \gamma p + \delta q$, patetque $F$ transire in $f$ positis $x = \alpha t - \beta u$, $y = \gamma t - \delta u$, hancque transformationem fore impropriam.

Quodsi igitur formae $F$, $f$ tam proprie quam improprie sunt aequivalentes, inveniri poterit tam transformatio propria aliqua quam impropria.

\bigskip

\textbf{179.}

\textbf{Problema.} Si formae $F$, $f$ sunt aequivalentes: invenire omnes transformationes formae $F$ in $f$.

\textit{Sol.} Si formae $F$, $f$ unico tantum modo sunt aequivalentes i.e.~proprio tantum vel improprie tantum: quaeratur per art.~prae.~transformatio una formae $F$ in $f$, patetque alias quam quae huic sint similes, dari non posse. Si vero formae $F$, $f$ tam proprie quam improprie aequivalent, quaerantur duae transformationes, altera propria, altera impropria. Iam sit forma $F = (A,B,C)$, $BB - AC = -D$, numerorumque $A$, $2B$, $C$ divisor communis maximus $= m$. Tum ex art.~162 patet, in priori casu omnes transformationes formae $F$ in $f$ ex una transformatione, in posteriori omnes proprias ex propria omnesque improprias ex impropria deduci posse, si modo omnes solutiones aequationis $tt + Duu = mm$ habeantur. His igitur inventis problema erit solutum.

Habetur autem $D = AC - BB$, $4D = 4AC - 4BB$, quare $\frac{4D}{mm} = 4\frac{A}{m}\frac{C}{m} - (\frac{2B}{m})^2$ erit integer. Iam si

1) $\frac{4D}{mm} > 4$, erit $D > mm$: quare in $tt + Duu = mm$, $u$ necessario debet esse $= 0$, adeoque $t$ alios valores quam $+m$ et $-m$ habere nequit. Hinc si $F$, $f$ unico tantum modo aequivalentes sunt et transformatio aliqua
\[
X = \alpha x' + \beta y', \quad Y = \gamma x' + \delta y'
\]
praeter hanc ipsam, quae prodit ex $t = m$ (art.~162), et hanc
\[
x = -\alpha x' - \beta y', \quad y = -\gamma x' - \delta y'
\]
aliae locum habere non possunt. Si vero $F$, $f$ tum proprie tum improprie aequivalent, atque propria aliqua transformatio habetur
\[
X = \alpha x' + \beta y', \quad Y = \gamma x' + \delta y'
\]
impropriaque
\[
X = \alpha' x' + \beta' y', \quad Y = \gamma' x' + \delta' y'
\]
praeter illam (ex $t = m$) hancce
\[
x = -\alpha x' - \beta y', \quad y = -\gamma x' - \delta y' \quad (\text{ex } t = -m)
\]
alia propria non dabitur; similiterque nulla impropria praeter
\[
x = -\alpha' x' - \beta' y', \quad y = -\gamma' x' - \delta' y' \quad (\text{ex } t = -m)
\]
et $x = \alpha' x' + \beta' y'$, $y = \gamma' x' + \delta' y'$ (ex $t = m$).

2) Si $\frac{4D}{mm} = 4$, sive $D = mm$, aequatio $tt + Duu = mm$ quatuor solutiones admittet: $t,u = m,0$; $-m,0$; $0,1$; $0,-1$. Hinc $F$, $f$ unico tantum modo aequivalentes sunt et transformatio aliqua
\[
X = \alpha x' + \beta y', \quad Y = \gamma x' + \delta y'
\]
quatuor omnino transformationes dabuntur,
\[
\begin{aligned}
x &= +\alpha x' + \beta y', & y &= +\gamma x' + \delta y';\\
x &= -\alpha x' - \beta y', & y &= -\gamma x' - \delta y';\\
x &= +\gamma x' - \delta y', & y &= -\alpha x' + \beta y';\\
x &= -\gamma x' + \delta y', & y &= +\alpha x' - \beta y'.
\end{aligned}
\]
Si vero $F$, $f$ duobus modis aequivalent, sive praeter transformationem illam datam alia ipsi dissimilis habetur, haec quoque suppeditabit quatuor illis dissimiles, ita ut octo transformationes habeantur. Ceterum facile demonstrari potest, in hoc casu $F$, $f$ semper revera duobus modis aequivalere. Nam quum $D = mm = AC - BB$, $m$ etiam ipsum $B$ metietur. Formae $(\frac{A}{m}, \frac{B}{m}, \frac{C}{m})$ determinans $= -1$, quare formae erit vel huic $(1,0,1)$ vel huic $(-1,0,-1)$ aequivalens.

Facile vero perspicitur, per eandem transformationem, per quam $(\frac{A}{m}, \frac{B}{m}, \frac{C}{m})$ transeat in $(+1,0,+1)$ formam $(A,B,C)$ transire in $(+m,0,+m)$, ancipitem. Quare forma $(A,B,C)$, ancipiti aequivalens, cuivis formae, cui aequivalet, tum proprie tum improprie aequivalebit.

3) Si $\frac{4D}{mm} = 3$, sive $4D = 3mm$. Tum $m$ erit par omnesque solutiones aequationis $tt + Duu = mm$ erunt sex,
\[
t,u = m,0; \; -m,0; \; \tfrac{1}{2}m, 1; \; -\tfrac{1}{2}m, -1; \; \tfrac{1}{2}m, -1; \; -\tfrac{1}{2}m, 1.
\]
Si itaque duae transformationes dissimiles formae $F$ in $f$ habentur
\[
x = \alpha x + \beta y, \quad y = \gamma x + \delta y
\]
habebuntur duodecim transformationes, scilicet sex priori similes
\[
x = \pm \alpha x' + \beta y', \quad y = \pm \gamma x' + \delta y'
\]
et sex posteriori similes, quae ex his nascuntur ponendo pro $\alpha$, $\beta$, $\gamma$, $\delta$ hos $\alpha'$, $\beta'$, $\gamma'$, $\delta'$.

Quod vero in hoc casu semper $F$, $f$ utroque modo aequivalent, ita demonstramus. Formae $(\frac{A}{m}, \frac{B}{m}, \frac{C}{m})$ determinans erit $= -\frac{3}{4}mm = -3$, adeoque haec forma (art.~176) aut formae $(+1,0,+3)$ aut huic $(+2,+1,+2)$ aequivalens. Unde facile perspicitur, formam $(A,B,C)$ aut formae $(+4m, 0, +\frac{1}{4}m)$ aut huic $(+m, +\frac{1}{2}m, +m)$\footnote{Demonstrari potest, formam $(A,B,C)$ necessario posteriori aequivalere: sed hoc hic non necessarium.}, quae ambae sunt ancipites, aequivalere adeoque, cuivis aequivalenti, utroque modo.

4) Si supponitur $\frac{4D}{mm} = 2$, fit $4D = 4mm - 2$, adeoque $= 2 \pmod{4}$. Sed quum nullum quadratum esse possit $= 2 \pmod{4}$, hic casus locum habere nequit.

5) Supponendo $\frac{4D}{mm} = 1$, fit $4D = 4mm - 1 \equiv -1 \pmod{4}$. Quod quum impossibile sit, etiam hic casus nequit locum habere.

Ceterum quum $D$ neque $= 0$, neque negativus sit, alii casus praeter enumeratos dari non possunt.

\bigskip

\textbf{180.}

\textbf{Problema.} Invenire omnes repraesentationes numeri dati $M$ per formam $axx + 2bxy + cyy\ldots F$ determinantis negativi $-D$, in quibus $x,y$ valores inter se primos nanciscuntur.

\textit{Sol.} Ex art.~154 patet, $M$ eo quo requiritur modo repraesentari non posse, nisi $-D$ sit resid.~quadr.~ipsius $M$. Investigentur itaque primo omnes valores diversi (i.e.~incongrui) expr.~$\sqrt{-D} \pmod{M}$, qui sint $N$, $-N$, $N'$, $-N'$, $N''$, $-N''$ etc.; quo simplicior evadat calculus, omnes $N$, $N'$ etc.~ita determinari possunt, ut non sint $> \frac{1}{2}M$. Iam quoniam quaevis repraesentatio ad aliquem horum valorum pertinere debet, singuli seorsim considerentur.

Si formae $F$, $(M, N, \frac{NN-D}{M})$ non sunt proprie aequivalentes, nulla repraesentatio ipsius $M$ ad valorem $N$ pertinens dari potest (art.~168). Si vero sunt, investigetur transformatio propria formae $F$ in $Mxx + 2Nxy + \frac{NN-D}{M}yy$, quae sit
\[
X = \alpha x + \beta y, \quad Y = \gamma x + \delta y;
\]
eritque $x = \alpha$, $y = \gamma$ repraesentatio numeri $M$ per $F$ ad $N$ pertinens. Sit div.~comm.~max.~numerorum $A$, $2B$, $C$, $= m$; distinguanturque tres casus (art.~prae.):

1) Si $\frac{4D}{mm} > 4$, aliae repraesentationes ad $N$ pertinentes quam hae duae $x = \pm \alpha$, $y = \pm \gamma$; $x = \mp \alpha$, $y = \mp \gamma$ non dabuntur (artt.~169,~179).

2) Si $\frac{4D}{mm} = 4$, habebuntur quatuor repraesentationes.

3) Si $\frac{4D}{mm} = 3$, habebuntur sex repraesentationes.

Eodem modo quaerendae sunt repraesentationes ad valores $-N$, $N'$, $-N'$ etc.~pertinentes.

\bigskip

\textbf{181.}

Investigatio repraesentationum numeri $M$ per formam $F$, in quibus $x$, $y$ valores inter se non primos habent, ad casum iam consideratum facile reduci potest. Fiat talis repraesentatio ponendo $x = \mu e$, $y = \mu f$, ita ut $\mu$ sit div.~comm.~max.~ipsorum $x,y$, i.e.~sive $e$, $f$ inter se primi. Tum erit $M = \mu\mu(Aee + 2Bef + Cff)$ adeoque per $\mu\mu$ divisibilis; substitutio vero $x = e$, $y = f$ erit repraesentatio numeri $\frac{M}{\mu\mu}$ per formam $F$, in qua $x$, $y$ valores inter se primos habent. Si itaque $M$ per nullum quadratum (praeter $1$) divisibilis est, e.g.~si est numerus primus, tales repraesentationes ipsius $M$ non dabuntur.

Si vero $M$ divisores quadraticos implicat, sint hi $\mu\mu$, $\nu\nu$, $\tau\tau$ etc. Quaerantur primo omnes repraesentationes numeri $\frac{M}{\mu\mu}$ per formam $(A,B,C)$, in quibus $x$, $y$ valores inter se primos habent, qui valores si per $\mu$ multiplicantur praebebunt omnes repraesentationes ipsius $M$, in quibus div.~comm.~max.~numerorum $x$, $y$ est $\mu$. Simili modo omnes repraesentationes ipsius $\frac{M}{\nu\nu}$, in quibus valores ipsorum $x$, $y$ inter se primi sunt, praebebunt omnes repraesentationes ipsius $M$, in quibus div.~comm.~max.~valorum ipsorum $x$, $y$ est $\nu$ etc.

Palam igitur est, per praecepta praecedentia omnes repraesentationes numeri dati per formam datam determinantis negativi inveniri posse.

\bigskip

\begin{center}
\Large Applicationes speciales ad discerptionem numerorum in quadrata duo,\\
in quadratum simplex et duplex, in simplex et triplex.
\end{center}

\bigskip

\textbf{182.}

Descendimus ad quosdam casus particulares, tum propter insignem ipsorum elegantiam tum propter assiduam operam ab ill.~Euler ipsis impensam, unde classicam quasi dignitatem sunt nacti.

\textbf{I.} Per formam $xx + yy$ ita repraesentari, ut $x$ ad $y$ sit primus (sive in duo quadrata inter se prima discerpi), nullus numerus potest nisi cuius residuum quadraticum est $-1$, tales vero numeri, positive accepti, omnes poterunt. Sit $M$ talis numerus, omnesque valores expr.~$\sqrt{-1} \pmod{M}$ hi: $N$, $-N$, $N'$, $-N'$, $N''$, $-N''$ etc. Tum per art.~176 forma $(M, N, \frac{NN+1}{M})$ formae $(1,0,1)$ proprie aequivalens erit. Sit transformatio aliqua propria huius in illam, $x = \alpha x' + \beta y'$, $y = \gamma x' + \delta y'$; eruntque repraesentationes numeri $M$ per formam $xx + yy$ ad $N$ pertinentes hae quatuor\footnote{Patet enim, hunc casum sub (2) art.~180 contentum esse.}: $x = +\alpha$, $y = +\gamma$; $x = -\gamma$, $y = +\alpha$; $x = +\gamma$, $y = -\alpha$; $x = -\alpha$, $y = -\gamma$.

Quum forma $(1,0,1)$ sit anceps, patet, etiam formam $(M,N,\frac{NN+1}{M})$ ipsi proprie aequivalentem fore, illamque proprie in hanc transmutari positis $X = \alpha x - \beta y$, $y = -\gamma x + \delta y$. Hinc derivantur quatuor repraesentationes ipsius $M$ ad $-N$ pertinentes, $x = +\alpha$, $y = -\gamma$; $x = +\gamma$, $y = +\alpha$; $x = -\gamma$, $y = -\alpha$; $x = -\alpha$, $y = +\gamma$. Manifestum itaque est, octo repraesentationes ipsius $M$ dari, quarum semissis altera ad $N$, altera ad $-N$ pertineat; sed hae omnes unicam tantummodo discerptionem numeri $M$ in duo quadrata exhibent, $M = \alpha\alpha + \gamma\gamma$, siquidem ad quadrata ipsa tantum neque vero ad ordinem radicumve signa spectamus.

Quodsi itaque alii valores expr.~$\sqrt{-1} \pmod{M}$ praeter $N$ et $-N$ non dantur, quod e.g.~evenit, quando $M$ est numerus primus, $M$ unico tantum modo in duo quadrata inter se prima resolvi poterit. Iam quum $-1$ sit residuum quadraticum cuiusvis numeri primi formae $4n+1$ (art.~108), manifestoque numerus primus in duo quadrata inter se non prima discerpi nequeat, habemus theorema

\begin{quote}
\textbf{Quivis numerus primus formae $4n+1$ in duo quadrata decomponi potest, et quidem unico tantum modo.}
\end{quote}

$2 = 1+1$, $5 = 1+4$, $13 = 4+9$, $17 = 1+16$, $29 = 4+25$, $37 = 1+36$, $41 = 16+25$, $53 = 4+49$, $61 = 25+36$, $73 = 9+64$, $89 = 25+64$, $97 = 16+81$ etc.

Theorema hoc elegantissimum iam Fermatio notum fuit, sed ab ill.~Euler primo demonstratum est, Comm.~nov.~Petr.~T.~V, ad annos 1754, 1755, p.~3 sqq. In T.~IV, diss.~exstat ad idem argumentum pertinens, p.~3 sqq., sed tum rem penitus nondum absolverat, vid.~imprimis art.~27.

Si igitur numerus aliquus formae $4n+1$ aut pluribus modis aut nullo modo in duo quadrata resolvi potest, certo non erit primus.

Vice versa autem, si expr.~$\sqrt{-1} \pmod{M}$ praeter $N$ et $-N$ alios adhuc valores habet, aliae adhuc repraesentationes ipsius $M$ dabuntur, ad hos pertinentes. In hoc itaque casu $M$ pluribus modis in duo quadrata resolvi poterit e.g.
\[
65 = 1 + 64 = 16 + 49, \quad 221 = 25 + 196 = 100 + 121.
\]

Repraesentationes reliquae, in quibus $x$, $y$ valores obtinent non primos inter se, per methodum nostram generalem facile inveniri possunt. Observamus tantummodo, si numerus aliquis factores formae $4n+3$ involvens, per nullam divisionem per quadratum ab his liberari possit (quod fiet, si aliquis aut plures talium factorum dimensionem imparem habet), hunc nullo modo in duo quadrata resolvi posse\footnote{Si numerus $M = 2^\mu \cdot 3^\alpha \cdot 5^\beta \cdot 7^\gamma \ldots$ ita ut $a$, $b$, $c$ etc.~sint numeri primi inaequales formae $4n+1$, atque $S$ productum ex omnibus factoribus primis ipsius $M$ formae $4n+3$ (ad quam formam quivis numerus positivus reduci potest, faciendo $S = 1$, quando $M$ est impar, et $S = 2$, quando $M$ nullos factores formae $4n+3$ implicat): $M$ nullo modo in duo quadrata resolvi poterit, si $S$ est non-quadratus; si vero $S$ est quadratus, dabuntur $\frac{1}{2}(\alpha+1)(\beta+1)(\gamma+1)\ldots$ discerptiones ipsius $M$, quando aliquis numerorum $\alpha$, $\beta$, $\gamma$ etc.~est impar, aut $\frac{1}{2}(\alpha+1)(\beta+1)(\gamma+1)\ldots + \frac{1}{2}$, quando omnes $\alpha$, $\beta$, $\gamma$ etc.~sunt pares (siquidem ad quadrata ipsa tantum respicitur). Qui in calculo combinationum aliquantum sunt versati, demonstrationem huius theorematis (cui, perinde ut aliis particularibus, immorari nobis non licet) ex theoria nostra generali haud difficulter eruere poterunt. Cf.~art.~105.}.

\textbf{II.} Per formam $xx + 2yy$ nullus numerus, cuius non-residuum $-2$, ita repraesentari potest, ut $x$ ad $y$ sit primus, reliqui omnes poterunt. Sit $-2$ residuum numeri $M$, atque $N$ valor aliquis expr.~$\sqrt{-2} \pmod{M}$. Tum per art.~176 formae $(1,0,2)$, $(M,N,\frac{NN+2}{M})$ proprio aequivalentes erunt. Transeat illa proprie in hanc ponendo $x = \alpha x' + \beta y'$, $y = \gamma x' + \delta y'$; eritque $x = \alpha$, $y = \gamma$ repraesentatio numeri $M$ ad $N$ pertinens. Praeter hanc et $x = -\alpha$, $y = -\gamma$ aliae ad $N$ non pertinebunt (art.~180).

Simili modo, ut supra, perspicitur, repraesentationes $x = \pm\alpha$, $y = \mp\gamma$ ad valorem $-N$ pertinere. Omnes vero hae quatuor repraesentationes unicam tantum discerptionem ipsius $M$ in quadratum et quadratum duplex exhibent, et si praeter $N$ et $-N$ alii valores expr.~$\sqrt{-2} \pmod{M}$ non dantur, aliae discerptiones non dabuntur. Hinc adiumento proposs.~art.~116 facile deducitur theorema

\begin{quote}
\textbf{Quivis numerus primus formae $8n+1$, vel $8n+3$, in quadratum et quadratum duplex decomponi potest et quidem unico tantum modo.}
\end{quote}

$1 = 1 + 0$, $3 = 1 + 2$, $11 = 9 + 2$, $17 = 9 + 8$, $19 = 1 + 18$, $41 = 9 + 32$, $43 = 25 + 18$, $59 = 9 + 50$, $67 = 49 + 18$, $73 = 1 + 72$, $83 = 81 + 2$, $89 = 81 + 8$, $97 = 25 + 72$ etc.

Etiam hoc theorema, ut plura similia, Fermatio innotuit: sed ill.~La Grange primus demonstrationem dedit, Suite des recherches d'Arithmetique, Nouv.~Mem.~de l'Ac.~de Berlin 1775, p.~323 sqq. Multa ad idem argumentum pertinentia iam ill.~Euler absolverat, Specimen de usu observationum in mathesi pura Comm.~nov.~Petr.~T.~VI, p.~185 sqq. Sed demonstratio completa theorematis semper ipsius industriam elusit, p.~220. Conf.~etiam diss.~in T.~VIII (ad annos 1760, 1761), Supplementum quorundam theorematum arithmeticorum, sub fin.

\textbf{III.} Per methodum similem demonstratur, quemvis numerum, cuius residuum quadr.~sit $-3$, repraesentari posse aut per formam $xx + 3yy$, aut per hanc $2xx + 2xy + 2yy$, ita ut valor ipsius $x$ ad valorem ipsius $y$ sit primus. Quare quum $-3$ sit residuum omnium numerorum primorum formae $3n+1$ (art.~119) manifestoque per formam $2xx + 2xy + 2yy$ numeri pares tantum repraesentari possint: eodem modo ut supra habetur theorema:

\begin{quote}
\textbf{Quivis numerus primus formae $3n+1$ in quadratum et quadratum triplex decomponi potest, et quidem unico tantum modo.}
\end{quote}

$1 = 1 + 0$, $7 = 4 + 3$, $13 = 1 + 12$, $19 = 16 + 3$, $31 = 4 + 27$, $37 = 25 + 12$, $43 = 16 + 27$, $61 = 49 + 12$, $67 = 64 + 3$, $73 = 25 + 48$ etc.

Demonstrationem huius theorematis ill.~Euler primus tradidit in commentatione modo laudata, Comm.~nov.~Petr.~T.~VIII, p.~105 sqq.

Simili modo ulterius progredi et e.g.~ostendere possemus, quemvis numerum primum formae $20n+1$, vel $20n+3$, vel $20n+7$, vel $20n+9$ (quippe quorum residuum $= -5$) per alterutram formam $xx + 5yy$, $2xx + 2xy + 3yy$ repraesentari posse, et quidem numeros primos formae $20n+1$ et $20n+9$ per priorem, primos formae $20n+3$, $20n+7$ per posteriorem, nec non dupla primorum formae $20n+1$, $20n+9$ per formam $2xx + 2xy + 3yy$, dupla primorum formae $20n+3$, $20n+7$ per formam $xx + 5yy$: sed hanc propositionem infinitasque alias particulares quivis proprio marte ex praecedentibus et infra tradendis derivare poterit. Transimus itaque ad formas determinantem positivi, et quum harum indoles prorsus alia sit, quando determinans est quadratus, alia, quando non-quadratus: formas determinantis quadrati hic primo excludimus posteaque seorsim considerabimus.


\newpage

\begin{center}
\Large De formis determinantis positivi non-quadrati.
\end{center}

\bigskip

\textbf{183.}

\textbf{Problema.} Proposita forma quacunque $(a,b,a')$, cuius determinans positivus non-quadratus $= D$: invenire formam huic proprie aequivalentem, $(A,B,C)$, in qua $B$ sit positivus et $< \sqrt{D}$; $A$ vero si est positivus, vel $-A$, si $A$ negativus, inter $\sqrt{D}+B$ et $\sqrt{D}-B$ situs.

\textit{Sol.} Supponimus in forma proposita utramque conditionem nondum locum habere; alioquin enim aliam formam quaerere opus non esset. Porro observamus, in forma determinantis non-quadrati terminum primum vel ultimum $= 0$ esse non posse (art.~171 ann.). Sit $b' \equiv -b \pmod{a'}$ atque intra limites $\sqrt{D}$ et $\sqrt{D} \mp a'$ situs (accepto signo superiore, quando $a$ positivus, inferiori, quando est negativus), quod fieri posse simili ratione ut art.~3, facile demonstratur; ponaturque $\frac{b'b' - D}{a'} = a''$, qui erit integer, quia $b'b' - D \equiv bb - D = aa' \pmod{a}$. Iam si $a'' < a'$, fiat denuo $b'' \equiv -b' \pmod{a''}$ et inter $\sqrt{D}$ et $\sqrt{D} \mp a''$ situs (prout $a'$ positivus vel negativus) et $\frac{b''b'' - D}{a''} = a'''$. Si hic iterum $a''' < a''$, sit rursus $b''' \equiv -b'' \pmod{a'''}$ et inter $\sqrt{D}$ et $\sqrt{D} \mp a'''$ situs atque $\frac{b'''b''' - D}{a'''} = a^{IV}$. Haec operatio continuetur, donec in progressione $a'$, $a''$, $a'''$, $a^{IV}$ etc.~ad terminum $a^m$ perveniatur, praecedente $a^{m-1}$ non minorem, quod tandem evenire debet, quia alioquin progressio infinita numerorum integrorum continuo decrescentium haberetur. Tum positis $a^m = A$, $b^m = B$, $a^{m+1} = C$, forma $(A,B,C)$ omnibus conditionibus satisfaciet.

\textit{Dem.} I. Quoniam in progressione formarum $(a,b,a')$, $(a',b',a'')$, $(a'',b'',a''')$ etc.~quaevis praecedenti est contigua: ultima $(A,B,C)$ primae $(a,b,a')$ proprie aequivalens erit.

II. Quia $B$ inter $\sqrt{D}$ et $\sqrt{D} \mp A$ situs est (accipiendo semper signum superius, quando $A$ est positivus, inferius, quando $A$ est negativus) patet, si ponatur $\sqrt{D} - B = p$, $B - (\sqrt{D} \mp A) = q$, hos $p$, $q$ fore positivos. Iam facile confirmatur, fore $qq + 2pq + 2p\sqrt{D} = D + AA - BB$; quare $D + AA - BB$ erit numerus positivus, quem ponemus $= r$. Hinc propter $D = BB - AC$, fit $r = AA - AC$, adeoque $AA - AC$ numerus positivus: quia vero per hyp. $A$ non est maior quam $C$, manifeste illud aliter fieri nequit, quam $AC$ est negativus, adeoque signa ipsorum $A$, $C$ opposita. Hinc $BB = D + AC < D$ adeoque $B < \sqrt{D}$.

III. Porro quia $-AC = D - BB$, erit $AC < D$, et hinc (quia $A$ non $> C$), $A < \sqrt{D}$. Quare $\sqrt{D} \mp A$ erit positivus, adeoque etiam $B$, qui inter limites $\sqrt{D}$ et $\sqrt{D} \mp A$ est situs.

IV. Hinc a potiori $\sqrt{D} + B > A$ positivus, et quia $\sqrt{D} - B > A = -q$, est negativus, $+A$ situs erit inter $\sqrt{D} + B$ et $\sqrt{D} - B$. \hfill Q.E.D.

\textit{Ex.} Proposita sit forma $(67, 97, 140)$, cuius determinans $= 29$. Hic invenitur progressio formarum $(67, 97, 140)$, $(140, -97, 67)$, $(67, -37, 20)$, $(20, -3, 1)$, $(-1, 5, 4)$. Ultima erit quaesita.

Tales formas $(A, B, C)$ determinantis positivi non-quadrati $D$, in quibus $A$ positive acceptus iacet inter $\sqrt{D} + B$ et $\sqrt{D} - B$, $B$ vero positivus est atque $< \sqrt{D}$, \textbf{formas reductas} vocabimus. Formae itaque reductae determinantis positivi non-quadrati aliquantum differunt a formis reductis determinantis negativi; sed propter magnam analogiam inter has et illas, denominationes diversas introducere noluimus.

\bigskip

\textbf{184.}

Si aequivalentia duarum formarum reductarum determinantis positivi aeque facile dignosci posset, ut in formis determinantis negativi (art.~172), aequivalentiam duarum formarum quarumcunque eiusdem determinantis positivi nullo negotio diiudicare possemus. Sed hic res longe aliter se habet, fierique potest ut permultae formae reductae inter se aequivalentes sint. Antequam itaque problema hoc aggrediamur, profundius in naturam formarum reductarum (determinantis positivi non-quadrati, quod semper hic subintelligendum) inquirere necesse erit.

1) Si $(a,b,c)$ est forma reducta, $a$ et $c$ signa opposita habebunt. Nam posito determinante formae $= D$, erit $ac = bb - D$ adeoque, propter $b < \sqrt{D}$, negativus.

2) Numerus $c$ perinde ut $a$, positive acceptus, inter $\sqrt{D} + b$ et $\sqrt{D} - b$ situs erit. Nam $-c = \frac{D - bb}{a}$; quare, abstractione facta a signo, $c$ iacebit inter $\sqrt{D} + b$ et $\sqrt{D} - b$.

3) Hinc patet, etiam $(c, b, a)$ fore formam reductam.

4) Tum $a$ tum $c$ erunt $< 2\sqrt{D}$. Uterque enim est $< \sqrt{D} + b$, adeoque a potiori $< 2\sqrt{D}$.

5) Numerus $b$ situs erit inter $\sqrt{D}$ et $\sqrt{D} \mp a$ (accepto signo superiori, quando $a$ positivus, inferiori, quando est negativus). Quia enim $+a$ iacet inter $\sqrt{D} + b$ et $\sqrt{D} - b$, erit $+a - (\sqrt{D} - b)$, sive $b - (\sqrt{D} \mp a)$ positivus; $b - \sqrt{D}$ autem est negativus; quamobrem $b$ inter $\sqrt{D}$ et $\sqrt{D} \mp a$ erit situs. Prorsus eodem modo demonstratur, $b$ inter $\sqrt{D}$ et $\sqrt{D} \mp c$ iacere (prout $c$ pos.~vel neg.).

6) Cuivis formae reductae $(a,b,c)$ ab utraque parte contigua est reducta una, et non plures.

Fiat $a' = c$, $b' \equiv -b \pmod{a'}$ et inter $\sqrt{D}$ et $\sqrt{D} \mp a'$ situs\footnote{Ubi signa ambigua sunt, superiora semper valent, quando $a'$ est positivus, inferiora, quando $a'$ negativus.}, $c' = \frac{bb - D}{a'}$; eritque forma $(a', b', c')$ formae $(a,b,c)$ ab ultima parte contigua, simulque manifestum est, si ulla forma reducta formae $(a,b,c)$ ab ultima parte contigua detur, eam ab hac $(a', b', c')$ diversam esse non posse. Hanc vero revera esse reductam, ita demonstramus.

A) Si ponitur
\[
\sqrt{D} + b = p, \quad +a - (\sqrt{D} - b) = q, \quad \sqrt{D} - b = r
\]
hi $p$, $q$, $r$ ex (2) supra et defin.~formae reductae erunt positivi. Porro ponatur
\[
b' - (\sqrt{D} \mp a') = q', \quad \sqrt{D} - b' = r'
\]
eruntque $q'$, $r'$ positivi, quia $b'$ iacet inter $\sqrt{D}$ et $\sqrt{D} \mp a'$. Denique sit $b + b' = \pm ma'$, eritque $m$ integer. Iam patet esse $p + q = b + b'$, adeoque $b + b'$ sive $\pm ma'$ positivum, et proin etiam $m$; unde sequitur, $m - 1$ certe non esse negativum. Porro fit
\[
r + q' + md = 2b' + d, \quad \text{sive } 2b' = r + q' + (m - 1)d
\]
umde $2b'$ et $b'$ necessario erunt positivi. Et quoniam $b' + r' = \sqrt{D}$, erit $b' < \sqrt{D}$.

B) Porro fit
\[
r + md = \sqrt{D} + b', \quad \text{sive } r + (m - 1)a' = \sqrt{D} - b' + a'
\]
quare $\sqrt{D} + b' > a'$ erit positivus. Hinc et quoniam $+a' - (\sqrt{D} - b') = q'$, adeoque positivus, $+a'$ iacebit inter $\sqrt{D} + b'$ et $\sqrt{D} - b'$. Quocirca $(a', b', c')$ erit forma reducta.

Eodem modo demonstratur, si fiat $'c = a$, $'b \equiv -b \pmod{'c}$ et inter $\sqrt{D}$ et $\sqrt{D} \mp 'c$ situs, $'a = \frac{bb - D}{'c}$: formam $('a, 'b, 'c)$ fore reductam. Manifesto autem forma haec formae $(a,b,c)$ a parte prima est contigua, aliaque reducta praeter $('a, 'b, 'c)$ hac proprietate praedita esse non poterit.

\textit{Ex.} Formae reductae $(5, 11, -14)$, cuius determinans $= 191$, a parte ultima contigua reducta $(-14, 3, 13)$, a parte prima vero haec $(-22, 9, 5)$.

7) Si formae reductae $(a,b,c)$ a parte ultima contigua est reducta $(a, b', c)$: reductae $(c, b, a)$ contigua erit a prima parte forma $(c, b', a)$; et si reductae $(a,b,c)$ a prima parte contigua est forma $('a, 'b, 'c)$: reductae $(c, b, a)$ reducta $(c, 'b, 'a)$ contigua erit ab ultima parte. Porro etiam formae $(-a, 'b, -'c)$, $(-a, b, -c)$, $(a', b', -c')$ reductae erunt, et secunda primae, tertia secundae ab ultima parte contiguae, sive prima secundae, secundaque tertiae a parte prima; similiterque tres formae $(-c, b', -a')$, $(-c, b, -a)$, $(-c, 'b, -'a)$. Haec tam obvia sunt, ut explicatione non egeant.

\bigskip

\textbf{185.}

Multitudo omnium formarum reductarum determinantis dati $D$ semper est finita, ipsae vero duplici modo inveniri possunt. Designemus indefinite omnes formas reductas determinantis $D$ per $(a,b,c)$, ita ut omnes valores ipsorum $a$, $b$, $c$ determinare oporteat.

\textbf{Methodus prima.} Accipiantur pro $a$ omnes numeri (tum positive, tum negative) minores quam $2\sqrt{D}$, quorum residuum quadraticum $= D$, et pro singulis $a$, ponatur $b$ aequalis omnibus valoribus positivis expr.~$\sqrt{D} \pmod{a}$ inter $\sqrt{D}$ et $\sqrt{D} \mp a$ iacentibus; $c$ vero pro singulis valoribus determinatis ipsorum $a$, $b$, ponatur $= \frac{bb - D}{a}$. Si quae formae hoc modo oriuntur, in quibus $+a$ extra $\sqrt{D} + b$ et $\sqrt{D} - b$ situs est, reiiciendae sunt.

\textbf{Methodus secunda.} Accipiantur pro $b$ omnes numeri positivi minores quam $\sqrt{D}$; pro singulis $bb - D$ resolvatur omnibus quibus fieri potest modis in binos factores, qui neglecto signo inter $\sqrt{D} + b$ et $\sqrt{D} - b$ iaceant, ponaturque alter $= a$, alter $= c$. Manifestum est, singulas resolutiones in factores praebere binas formas, quia uterque factor tum $= a$, tum $= c$ poni debet.

\textit{Ex.} Sit $D = 79$ eruntque valores ipsius $a$ viginti duo $\pm 1$, $2$, $3$, $5$, $6$, $7$, $9$, $10$, $13$, $14$, $15$. Unde inveniuntur formae undeviginti:
\[
\begin{aligned}
&(1,8,-15), &&(2,7,-15), &&(3,8,-5), &&(3,7,-10), &&(5,8,-3), &&(5,7,-6),\\
&(6,7,-5), &&(6,5,-9), &&(7,4,-9), &&(7,3,-10), &&(9,5,-6), &&(9,4,-7),\\
&(10,7,-3), &&(10,3,-7), &&(13,1,-6), &&(14,3,-5), &&(15,8,-1), &&(15,7,-2), &&(15,2,-5)
\end{aligned}
\]
totidemque aliae quae fiunt ex his, si terminorum exterorum signa commutantur, puta $(-1, 8, 15)$, $(-2, 7, 15)$ etc., ita ut omnes triginta octo sint. Sed ex his reiiciendae sex $(+13, 1, \mp 6)$, $(+14, 3, \mp 5)$, $(+15, 2, \mp 5)$, reliquae triginta duae omnes reductas amplectuntur. Per methodum secundam eaedem formae prodeunt sequenti ordine\footnote{Pro $b = 1$, $-78$ in duos factores, qui neglecto signo inter $\sqrt{79} + 1$ et $\sqrt{79} - 1$ iaceant, resolvi nequit; quare hic valor est praetereundus, ex eademque ratione valores $2$ et $6$.}:
\[
\begin{aligned}
&(+7,3,\mp 10), \; (+10,3,\mp 7), \; (+7,4,\mp 9), \; (+9,4,\mp 7), \; (+6,5,\mp 9), \; (+9,5,\mp 6),\\
&(+2,7,\mp 15), \; (+3,7,\mp 10), \; (+5,7,\mp 6), \; (+6,7,\mp 5), \; (+10,7,\mp 3), \; (+15,7,\mp 2),\\
&(+1,8,\mp 15), \; (+3,8,\mp 5), \; (+5,8,\mp 3), \; (+15,8,\mp 1).
\end{aligned}
\]

\bigskip

\textbf{186.}

Sit $F$ forma reducta determinantis $D$, ipsique ab ultima parte contigua forma reducta $F'$; huic iterum ab ultima parte contigua reducta $F''$; reducta $F'''$ ipsi $F''$ contigua ab ultima parte etc. Tum patet, omnes formas $F'$, $F''$, $F'''$ etc.~esse prorsus determinatas, et tum inter se tum formae $F$ proprie aequivalentes. Quoniam vero multitudo omnium formarum reductarum determinantis dati est finita, manifestum est, omnes formas in progressione infinita $F$, $F'$, $F''$ etc.~diversas esse non posse. Ponamus $F^m$ et $F^{m+n}$ esse identicas, eruntque $F^{m-1}$ et $F^{m+n-1}$ ei~dem formae reductae a parte prima contiguae, adeoque identicae; hinc eodem modo $F^{m-2}$ et $F^{m+n-2}$ etc.~tandemque $F$ et $F^n$ identicae erunt. Quare in progressione $F$, $F'$, $F''$ etc., si modo satis longe continuatur, necessario tandem forma prima $F$ recurret; et si supponimus $F^n$ esse primam identicam cum $F$, sive omnes $F'$, $F''$\ldots$F^{n-1}$ a forma $F$ diversas: facile perspicitur, omnes formas $F$, $F'$, $F''$\ldots$F^{n-1}$ diversas fore. Complexum harum formarum vocabimus \textbf{periodum} formae $F$. Si igitur progressio ultra ultimam periodi formam producitur, eaedem formae $F$, $F'$, $F''$ etc.~iterum prodibunt, progressioque tota infinita $F$, $F'$, $F''$ etc.~constituta erit ex hac periodo formae $F$ infinities repetita.

Progressio $F$, $F'$, $F''$ etc.~etiam retro continuari potest, praeponendo formae $F$ reductam $'F$, quae ipsi a parte prima est contigua; huic iterum reductam $''F$, quae ipsi a prima parte contigua etc. Hoc modo habebitur progressio formarum utrimque infinita
\[
\ldots \; '''F, \; ''F, \; 'F, \; F, \; F', \; F'', \; F''' \; \ldots
\]
perspicieturque facile, $'F$ identicam fore cum $F^{n-1}$, $''F$ cum $F^{n-2}$ etc.~adeoque progressionem etiam a laeva parte e periodo formae $F$, infinities repetita, esse constitutam.

Si formis $F$, $F'$, $F''$ etc.; $'F$, $''F$ etc.~tribuuntur indices $0$, $1$, $2$ etc.; $-1$, $-2$ etc., generaliterque formae $F^m$ index $m$, formae $^mF$ index $-m$: patet, formas quascunque seriei identicas fore vel diversas, prout ipsarum indices congrui sint vel incongrui secundum modulum $n$.

\textit{Ex.} Periodus formae $(3, 8, -5)$, cuius determinans $= 79$, invenitur haec: $(3, 8, -5)$, $(-5, 7, 6)$, $(6, 5, -9)$, $(-9, 4, 7)$, $(7, 3, -10)$, $(-10, 7, 3)$. Post ultimam iterum prodit $(3, 8, -5)$. Hic itaque $n = 6$.

\bigskip

\textbf{187.}

Ecce quasdam observationes generales circa has periodos.

1) Si formae $F$, $F'$, $F''$ etc.; $'F$, $''F$, $'''F$ etc.~ita exhibentur:
\[
(a, b, -a'), \; (-a', b', a''), \; (a'', b'', -a''') \; \ldots; \; (-a, 'b, a'), \; ('a, ''b, -''a) \; \ldots
\]
omnes $a$, $a'$, $a''$, $a'''$ etc.~$'a$, $''a$, $'''a$ etc.~eadem signa habebunt (art.~184,~1), omnes vero $b$, $b'$, $b''$ etc.~$'b$, $''b$ etc.~erunt positivi.

2) Hinc manifestum est, numerum $n$ (multitudinem formarum, ex quibus periodus formae $F$ constat) semper esse parem. Etenim terminus primus formae cuiusvis $F^m$ ex hac periodo manifesto idem signum habebit uti terminus primus formae $F$, si $m$ est par, oppositum, si $m$ est impar. Quare quum $F^n$ et $F$ identicae sint, $n$ necessario erit par.

3) Demonstravimus (art.~184), quomodo forma reducta $F^m$ periodi formae $F$ contigua sit a parte prima, nempe $= (a^m, b^{m-1}, -a^{m-1})$; et a parte posteriori, $= (a^m, b^{m+1}, -a^{m+1})$. Hinc facile colligitur, omnes periodos formarum reductarum determinantis dati $D$, quae quidem ex eodem termino $a$ incipiunt, eandem seriem terminorum $b$, $b'$, $b''$ etc.~habere.

4) Omnes formae reductae eiusdem periodi eandem signorum dispositionem habent. Signorum autem dispositionem vocabimus eam relationem, qua termini primi formarum $F$, $F'$, $F''$\ldots$F^{n-1}$ ad signa sua respiciunt. Ita in exemplo nostro signa terminorum primorum sunt $+$, $-$, $+$, $-$, $+$, $-$.

5) Formae reductae, quae diversis periodis pertinentes eandem signorum dispositionem habent, periodis suis in ordinem similem redactis, aequivalentes erunt. Ut si periodus formae $F$ redacta in ordinem est $(a, b, -a')$, $(-a', b', a'')$, $(a'', b'', -a''')$ etc., et periodus formae $G$ pariter est $(c, d, -c')$, $(-c', d', c'')$, $(c'', d'', -c''')$ etc.: formae $F$ et $G$ aequivalentes erunt. Hinc facile diiudicari poterit, an duae formae reductae eidem periodo an diversis adnumerandae sint.

6) Vocamus formas socias, quae ex iisdem terminis constant, sed ordine inverso positis, ut $(a, b, -a')$ et $(-a', b, a)$. Tum facile perspicitur ex art.~184,~7, si periodus formae reductae $F$ sit $F$, $F'$, $F''$\ldots$F^{n-1}$, formae $F$ socia $f$ formisque $F^{n-1}$, $F^{n-2}$\ldots$F'$, $F$ resp.~sociae sint formae $f'$, $f''$\ldots$f^{n-2}$, $f^{n-1}$: periodum formae $f$ fore $f$, $f'$, $f''$\ldots$f^{n-2}$, $f^{n-1}$ adeoque ex totidem formis constare, ut periodum formae $F$. Periodos formarum sociarum vocabimus periodos socias. Ita in exemplo nostro sociae sunt periodi III et VI; IV et V.

7) Sed fieri etiam potest, ut forma $f$ ipsa in periodo sociae suae $F$ occurrat, uti in ex.~nostro in periodo I et II, adeoque periodus formae $F$ cum periodo formae $f$ conveniat, sive ut periodus formae $F$ sibi ipsi sit socia. Quoties hoc evenit, in hac periodo duae formae ancipites invenientur. Ponamus enim periodum formae $F$ constare e $2m$ formis sive $F$ et $F^{2m}$ esse identicas; porro sit $m+1$ index formae $f$ in periodo formae $F$\footnote{Index hic necessario erit impar, quia manifesto termini primi formarum $F$, $f$ signa opposita habent (vid.~supra, 2).}, sive $F^{m+1} = f$ sociae. Tum patet etiam $F'$ et $F^{m+2}$ fore socias nec non $F''$ et $F^{m+3}$, unde etiam $F^m$ et $F^{2m+1}$. Sit $F^m = (a^m, b^m, -a^{m+1})$, $F^{m+1} = (-a^{m+1}, b^{m+1}, a^{m+2})$. Tum erit $b^m + b^{m+1} \equiv 0 \pmod{a^{m+1}}$; ex defin.~formarum sociarum vero erit $b^m = b^{m+1}$ atque hinc $2b^{m+1} \equiv 0 \pmod{a^{m+1}}$, sive forma $F^{m+1}$ anceps. Eodem modo $F^{2m+1}$ et $F^{2m}$ erunt sociae; hinc $F^{m-1}$ et $F^{m+2}$, $F^{m-2}$ et $F^{m+3}$ etc.~tandemque $F^{2m+1}$ et $F^{m}$, quarum posterior erit anceps, uti per simile ratiocinium facile probatur. Quia vero $m+1$ et $m+m+1$ secundum mod.~$m$ sunt incongrui, formae $F^{m+1}$ et $F^{m+m+1}$ identicae non erunt (art.~186, ubi $n$ idem denotat, quod hic $2m$). Ita in I sunt formae ancipites $(1, 8, -15)$, $(2, 7, -15)$; in II vero $(-1, 8, 15)$, $(-2, 7, 15)$.

8) Vice versa, quaevis periodus, in qua forma anceps occurrit, sibi ipsi socia est. Facile enim perspicitur, si $F^m$ sit forma reducta anceps: formam ipsi sociam (quae etiam est reducta) simul ipsi a parte prima contiguam esse, i.e.~$F^{m-1}$ et $F^m$ socias. Tum vero tota periodus sibi ipsi socia erit. Hinc patet, fieri non posse, ut unica tantum forma anceps in periodo aliqua contenta sit.

9) Sed etiam plures quam duae in eadem periodo esse nequeunt. Ponamus enim in periodo formae $F$, ex $2n$ formis constante, tres formas ancipites dari $F^\lambda$, $F^\mu$, $F^\nu$, ad indices $\lambda$, $\mu$, $\nu$ respective pertinentes, ita ut $\lambda$, $\mu$, $\nu$ sint numeri inaequales inter limites $0$ et $2n - 1$ (incl.) siti. Tum formae $F^{\lambda-1}$ et $F^\lambda$ erunt sociae; similiterque $F^{\mu-1}$ et $F^\mu$ etc.~tandemque $F^{\nu-1}$ et $F^\nu$. Ex eadem ratione $F^\lambda$ et $F^{2\lambda-1}$ sociae erunt, nec non $F^\mu$ et $F^{2\mu-1}$; quare $F^{2\lambda-1}$, $F^{2\mu-1}$, $F^{2\nu-1}$ identicae, indicesque $2\lambda - 1$, $2\mu - 1$, $2\nu - 1$ secundum modulum $2n$ congrui erunt, et proin etiam $\lambda \equiv \mu \equiv \nu \pmod{n}$. Q.E.A. quia manifeste inter limites $0$ et $2n - 1$ tres numeri diversi secundum modulum $n$ congrui iacere nequeunt.

\bigskip

\textbf{188.}

Quum omnes formae ex eadem periodo proprio sint aequivalentes, quaestio oritur, annon etiam formae e periodis diversis proprio aequivalentes esse possint. Sed antequam ostendamus, hoc esse impossibile, quaedam de transformatione formarum reductarum sunt exponenda.

Quoniam in sequentibus de formarum transformationibus persaepe agendum erit: ut prolixitatem quantum fieri potest evitemus, sequenti scribendi compendio abhinc semper utemur. Si forma aliqua $L'XX + 2M'XY + N'YY$ per substitutionem $X = \alpha x + \beta y$, $Y = \gamma x + \delta y$ in formam $lxx + 2mxy + nyy$ transformetur: simpliciter dicemus, $(L, M, N)$ transformari in $(l, m, n)$ per substitutionem $\alpha$, $\beta$, $\gamma$, $\delta$. Hoc modo opus non erit, indeterminatas formarum singularum, de quibus agitur, per signa propria denotare. Palam vero est, indeterminatam primam a secunda in quavis forma probe distingui debere.

Proposita sit forma reducta $(a, b, -a')\ldots f$, determinantis $D$. Formetur simili modo ut in art.~186 progressio formarum reductarum utrimque infinita,
\[
\ldots \; '''f, \; ''f, \; 'f, \; f, \; f', \; f'', \; f''' \; \ldots
\]
et quidem sit
\[
\begin{aligned}
f &= (a, b, -a'), & f' &= (-a', b', a''), & f'' &= (a'', b'', -a''') \; \text{etc.};\\
'f &= (-a, 'b, a'), & ''f &= ('a, ''b, -'a) \; \text{etc.}
\end{aligned}
\]
Ponatur
\[
\frac{b + b'}{-a'} = h', \quad \frac{b' + b''}{a''} = h'', \quad \frac{b'' + b'''}{-a'''} = h''' \; \text{etc.},
\]
\[
\frac{'b + b}{a} = 'h, \quad \frac{''b + 'b}{-'a} = ''h, \quad \frac{'''b + ''b}{''a} = '''h \; \text{etc.}
\]

Tum patet, si (ut in art.~177) numeri $\alpha'$, $\alpha''$, $\alpha'''$ etc.; $\beta'$, $\beta''$, $\beta'''$ etc.~etc.~formantur secundum algorithmum sequentem:
\[
\begin{aligned}
\alpha' &= 0, & \beta' &= -1, & \gamma' &= 1, & \delta' &= h';\\
\alpha'' &= \gamma', & \beta'' &= h''\gamma' - \beta', & \gamma'' &= \alpha', & \delta'' &= h''\alpha' - \delta';\\
\alpha''' &= \gamma'', & \beta''' &= h'''\gamma'' - \beta'', & \gamma''' &= \alpha'', & \delta''' &= h'''\alpha'' - \delta'' \; \text{etc.}
\end{aligned}
\]
$f$ transformatum iri
\begin{align*}
\text{in } &f' \text{ per substitutionem } \alpha', \beta', \gamma', \delta'\\
\text{in } &f'' \text{ per substitutionem } \alpha'', \beta'', \gamma'', \delta''\\
\text{in } &f''' \text{ per substitutionem } \alpha''', \beta''', \gamma''', \delta''' \text{ etc.}
\end{align*}
omnesque has transformationes fore proprias.

Quum $'f$ transeat in $f$ per substitutionem propriam $0$, $-1$, $1$, $h$ (art.~158): $f$ transibit in $'f$ per subst.~prop.~$h$, $1$, $-1$, $0$. Ex simili ratione $'f$ transibit in $''f$ per subst.~propr.~$'h$, $1$, $-1$, $0$; $''f$ in $'''f$ per subst.~pr.~$''h$, $1$, $-1$, $0$ etc. Hinc per art.~159 eodem modo ut art.~177 colligitur, si numeri $\alpha$, $''\alpha$, $'''\alpha$ etc.; $'\beta$, $''\beta$, $'''\beta$ etc.~etc.~formantur secundum algorithmum sequentem:
\[
\begin{aligned}
'\alpha &= h, & '\beta &= 1, & '\gamma &= -1, & '\delta &= 0;\\
''\alpha &= '\gamma, & ''\beta &= 'h'\gamma - '\beta, & ''\gamma &= '\alpha, & ''\delta &= 'h'\alpha - '\delta;\\
'''\alpha &= ''\gamma, & '''\beta &= ''h''\gamma - ''\beta, & '''\gamma &= ''\alpha, & '''\delta &= ''h''\alpha - ''\delta \; \text{etc.}
\end{aligned}
\]
$f$ transformatum iri
\begin{align*}
\text{in } &'f \text{ per substitutionem } '\alpha, '\beta, '\gamma, '\delta\\
\text{in } &''f \text{ per substitutionem } ''\alpha, ''\beta, ''\gamma, ''\delta\\
\text{in } &'''f \text{ per substitutionem } '''\alpha, '''\beta, '''\gamma, '''\delta\\
&\text{etc.}
\end{align*}
omnesque has transformationes fore proprias.

Si ponitur $\alpha = 1$, $\beta = 0$, $\gamma = 0$, $\delta = 1$: hi numeri eandem rationem habebunt ad formam $f$, quam habent $\alpha'$, $\beta'$, $\gamma'$, $\delta'$ ad $f'$; $\alpha''$, $\beta''$, $\gamma''$, $\delta''$ ad $f''$ etc.; $'\alpha$, $'\beta$, $'\gamma$, $'\delta$ ad $'f$ etc. Scilicet per substitutionem $\alpha$, $\beta$, $\gamma$, $\delta$ forma $f$ transibit in $f'$. Tum vero progressiones infinitae $\alpha$, $\alpha''$, $\alpha'''$ etc., $'\alpha$, $''\alpha$, $'''\alpha$ etc., per intercalationem termini $\alpha' = 0$, concinne iungentur, ita ut unam continuam utrimque infinitam constituere concipi possint secundum eandem legem ubique progredientem$\ldots'''\alpha$, $''\alpha$, $'\alpha$, $\alpha$, $\alpha'$, $\alpha''$, $\alpha'''$\ldots Lex progressionis haec est:
\[
'''\alpha + '\alpha = '''h \cdot ''\alpha, \quad ''\alpha + \alpha = ''h \cdot '\alpha, \quad '\alpha + \alpha' = 'h \cdot \alpha, \quad \alpha + \alpha'' = h' \cdot \alpha', \quad \alpha' + \alpha''' = h'' \cdot \alpha'' \quad \text{etc.}
\]
sive generaliter (si indicem negativum a dextra scriptum idem designare supponimus, ac positivum a laeva)
\[
\alpha^{m-1} + \alpha^{m+1} = h^m \alpha^m.
\]
Simili modo progressio $\ldots''\beta$, $'\beta$, $\beta$, $\beta'$, $\beta''$\ldots continua erit, cuius lex
\[
\beta^{m-1} + \beta^{m+1} = h^m \beta^m
\]
et proprio cum praecedente identica, omnibus terminis uno loco promotis, $''\beta = '\alpha$, $'\beta = \alpha$, $\beta = \alpha'$ etc. Lex progressionis continuae $\ldots''\gamma$, $'\gamma$, $\gamma$, $\gamma'$, $\gamma''$\ldots erit haec
\[
\gamma^{m-1} + \gamma^{m+1} = h^m \gamma^m
\]
et lex huius $\ldots''\delta$, $'\delta$, $\delta$, $\delta'$, $\delta''$\ldots erit
\[
\delta^{m-1} + \delta^{m+1} = h^m \delta^m.
\]
insuperque generaliter $h^m = \gamma^{m+1} \delta^m - \gamma^m \delta^{m+1}$.

\textit{Ex.} Sit forma proposita $f$ haec $(3, 8, -5)$, quae transformabitur in formam $'''f = (-10, 7, 3)$ per substitutionem $-805$, $-152$, $+143$, $+27$
\[
\begin{aligned}
'''f &= (-10, 7, 3) & \ldots & \quad -805, \; -152, \; +143, \; +27\\
''f &= (3, 8, -5) & \ldots & \quad -152, \; -27, \; +27, \; +5\\
'f &= (-5, 7, 6) & \ldots & \quad +27, \; +5, \; -10, \; -2\\
f &= (6, 5, -9) & \ldots & \quad -10, \; -2, \; +3, \; +1\\
f' &= (-9, 4, 7) & \ldots & \quad +3, \; +1, \; -1, \; 0\\
f'' &= (7, 3, -10) & \ldots & \quad -1, \; 0, \; 0, \; 1\\
f''' &= (-10, 7, 3) & \ldots & \quad 0, \; 1, \; 1, \; 0\\
f^{IV} &= (3, 8, -5) & \ldots & \quad 1, \; 0, \; 0, \; 1\\
f^V &= (-5, 7, 6) & \ldots & \quad 0, \; 1, \; 1, \; 0\\
f^{VI} &= (6, 5, -9) & \ldots & \quad 1, \; 0, \; 0, \; 1\\
f^{VII} &= (-9, 4, 7) & \ldots & \quad 0, \; 1, \; 1, \; 0 \quad \text{etc.}
\end{aligned}
\]

\bigskip

\textbf{189.}

1) Ex legibus progressionis in art.~prae.~traditis sequitur, omnes numeros $\alpha^m$, $\beta^m$, $\gamma^m$, $\delta^m$ (pro indice quocunque $m$) esse integros, quorum priores duo secundis ita sunt coniuncti, ut sit $\alpha^m \delta^m - \beta^m \gamma^m = (-1)^m$.

2) Hinc facile deducitur, fractionem $\frac{\gamma^m}{\alpha^m}$ ei~dem signo gaudere atque $a$ vel $a'$, prout $m$ est par vel impar, adeoque $\frac{\gamma^m}{\alpha^m}$ semper signum habere oppositum signo ipsius $a^{m+1}$; fractionem vero $\frac{\delta^m}{\beta^m}$ (quando denominator non est $= 0$) signum oppositum habere signo ipsius $a^m$, ita ut progressio $\frac{\gamma'}{\alpha'}$, $\frac{\gamma''}{\alpha''}$, $\frac{\gamma'''}{\alpha'''}$ etc.~continuo crescat, signumque termini $\alpha^m$ sit $+$, $-$, $-$, $+$, prout $m \equiv 0$, $1$, $2$, $3 \pmod{4}$.

3) Hoc modo invenitur, omnes quatuor progressiones infinitas $\alpha$, $\alpha'$, $\alpha''$ etc.; $\gamma$, $\gamma'$, $\gamma''$ etc.; $'\alpha$, $''\alpha$, $'''\alpha$ etc.; $'\gamma$, $''\gamma$, $'''\gamma$ etc.~continuo crescere, adeoque etiam sequentes cum illis identicas: $\beta$, $\beta'$, $\beta''$ etc.; $\delta$, $\delta'$, $\delta''$ etc.; $\beta'$, $\delta'$, $\beta''$, $\delta''$ etc.; $'\beta$, $''\beta$, $'''\beta$ etc.; et, prout $m \equiv 0$, $1$, $2$, $3 \pmod{4}$, signum
\[
\begin{array}{c|c|c}
\text{ipsius } \alpha^m, & + \quad - \quad - \quad + & \text{ipsius } \beta^m, \quad + \quad - \quad + \quad - \\
\text{ipsius } \gamma^m, & + \quad + \quad - \quad - & \text{ipsius } \delta^m, \quad + \quad + \quad - \quad - \\
\text{ipsius } {}^m\alpha, & + \quad + \quad - \quad - & \text{ipsius } {}^m\beta, \quad + \quad + \quad - \quad - \\
\text{ipsius } {}^m\gamma, & + \quad - \quad + \quad - & \text{ipsius } {}^m\delta, \quad + \quad - \quad + \quad -
\end{array}
\]
valentibus superioribus, quando $a$ est positivus, inferioribus, quando $a$ negativus.

Teneatur imprimis haec proprietas: Designante $m$ indicem quemcunque positivum, $\alpha^m$ et $\gamma^m$ habebunt eadem signa, quando $a$ positivus, opposita, quando $a$ negativus, similiterque $\beta^m$ et $\delta^m$; contra $'^m\alpha$ et $'^m\gamma$, vel $'^m\beta$ et $'^m\delta$ habebunt eadem signa, quando $a$ negativus, opposita, quando $a$ positivus.

4) In signis art.~27 magnitudo ipsorum $\alpha^m$ etc.~concinne ita exhiberi potest, ponendo
\[
\pm h' = k', \quad \pm h'' = k'', \quad \pm h''' = k''' \text{ etc.}, \quad \pm 'h = 'k, \quad \pm ''h = ''k \text{ etc.}
\]
ita ut omnes $k'$, $k''$ etc.~$'k$, $''k$ etc.~sint numeri positivi:
\[
\begin{aligned}
\alpha^m &= \pm [k', k'', k''' \ldots k^{m-1}], & \beta^m &= \pm [k', k'', \ldots k^m]\\
\gamma^m &= \pm [k'', k''' \ldots k^{m-1}], & \delta^m &= \pm [k', k'', k''' \ldots k^m]
\end{aligned}
\]
\[
\begin{aligned}
'^m\alpha &= \pm ['k, ''k \ldots {}^{m-1}k], & '^m\beta &= \pm ['k, ''k \ldots {}^mk]\\
'^m\gamma &= \pm [''k \ldots {}^{m-1}k], & '^m\delta &= \pm ['k, ''k \ldots {}^{m-1}k]
\end{aligned}
\]
signa vero ad praecepta modo tradita determinari debent. Secundum has formulas, quarum demonstrationem propter facilitatem omittimus, calculus semper expeditissime absolvi poterit.

\bigskip

\textbf{190.}

\textbf{Lemma.} Designantibus $m$, $\mu$, $m'$, $n$, $\nu$, $n'$ numeros integros quoscunque, ita tamen ut trium posteriorum nullus sit $= 0$: dico, si $\frac{\mu}{\nu}$ iaceat inter limites $\frac{m}{n}$ et $\frac{m'}{n'}$ excluive, atque sit $mn' - nm' = +1$, denominatorem $\nu$ fore maiorem quam $n$ et $n'$.

\textit{Dem.} Manifesto $\frac{\mu}{\nu} \cdot n n'$ iacebit inter $\frac{m}{n} \cdot n n'$ et $\frac{m'}{n'} \cdot n n'$; adeoque ab utroque limite minus differet quam limes alter ab altero, i.e.~erit $\frac{m}{n} n n' - \frac{\mu}{\nu} n n' < n n' \cdot |\frac{m}{n} - \frac{m'}{n'}|$ et $\frac{\mu}{\nu} n n' - \frac{m'}{n'} n n' < n n' \cdot |\frac{m}{n} - \frac{m'}{n'}|$, sive $\mu n' - \nu m < \frac{n'}{\nu}$ et $\nu m' - \mu n < \frac{n}{\nu}$. Hinc sequitur, quoniam $\mu n' - \nu m$ certe non $= 0$ (alioquin enim foret $\frac{\mu}{\nu} = \frac{m}{n'}$ contra hyp.), neque $\nu m' - \mu n = 0$ (ex simili ratione) sed uterque ad minimum $= 1$, fore $\nu > n'$ et $> n$. \hfill Q.E.D.

Perspicuum itaque est, $\nu$ non posse esse $= 1$, i.e.~si fuerit $mn' - nm' = +1$, inter fractiones $\frac{m}{n}$, $\frac{m'}{n'}$ nullum numerum integrum iacere posse. Quare etiam cifra inter ipsas iacere nequit, i.e.~fractiones istae signa opposita habere nequeunt.

\bigskip

\textbf{191.}

\textbf{Theorem.} Si forma reducta $(a, b, -a')$ determinantis $D$ per substitutionem $\alpha$, $\beta$, $\gamma$, $\delta$ transit in reductam $(A, B, -A')$ eiusdem determinantis: iacebit, primo, $\frac{-\gamma}{\alpha}$ inter $-\frac{\sqrt{D} + b}{a}$ et $-\frac{\sqrt{D} - b}{a}$ (siquidem neque $\gamma$ neque $\alpha = 0$, i.e.~si uterque limes est finitus), accepto signo superiore, quando neuter horum limitum habet signum signo ipsius $a$ oppositum (sive clarius, quando aut uterque idem habet, aut alter idem, alter est $= 0$), inferiori, quando neuter habet idem ut $a$; secundo $\frac{\delta}{\beta}$ inter $-\frac{\sqrt{D} + b}{a'}$ et $-\frac{\sqrt{D} - b}{a'}$ (siquidem neque $\alpha$ neque $\beta = 0$), signo superiori accepto, quando limes neuter signum signo ipsius $a'$ (vel $a$) oppositum habet, inferiori, quando neuter habet idem ut $a'$\footnote{Manifestum est, alios casus locum habere non posse, quum ex art.~prae.~propter $\alpha\delta - \beta\gamma = \pm 1$, limites bini neque signa opposita habere, neque simul $= 0$ esse possint.}.

\textit{Dem.} Habentur aequationes
\begin{align*}
a\alpha\alpha + 2b\alpha\gamma - a'\gamma\gamma &= A \tag{1}\\
a\beta\beta + 2b\beta\delta - a'\delta\delta &= -A' \tag{2}\\
\frac{\sqrt{D} + b}{a} \cdot \alpha + \frac{-\sqrt{D} + b}{a} \cdot \gamma &= \frac{-\sqrt{D} + B}{A} \text{ sive } \frac{\sqrt{D} + B}{A} \tag{3}\\
\frac{\sqrt{D} + b}{a} \cdot \beta + \frac{-\sqrt{D} + b}{a} \cdot \delta &= \frac{-\sqrt{D} - B}{-A'} \text{ sive } \frac{\sqrt{D} - B}{-A'} \tag{4}\\
\frac{\sqrt{D} + b}{a'} \cdot \gamma + \frac{-\sqrt{D} + b}{a'} \cdot \alpha &= \frac{\sqrt{D} + B}{A} \text{ sive } \frac{-\sqrt{D} + B}{A} \tag{5}\\
\frac{\sqrt{D} + b}{a'} \cdot \delta + \frac{-\sqrt{D} + b}{a'} \cdot \beta &= \frac{\sqrt{D} - B}{-A'} \text{ sive } \frac{-\sqrt{D} - B}{-A'} \tag{6}
\end{align*}

Aequatio 3, 4, 5, 6 reiicienda erit, si $\gamma$, $\delta$, $\alpha$, $\beta$ resp.~$= 0$. Sed dubium hic manet, quae signa quantitatibus radicalibus tribui debeant; hoc sequenti modo decidemus.

Statim patet, in (3) et (4) necessario signa superiora accipi debere, quando neque $-\frac{\sqrt{D} + b}{a}$ neque $-\frac{\sqrt{D} - b}{a}$ signum habeat signo ipsius $a$ oppositum; quoniam accepto signo inferiori $-\frac{\sqrt{D} + b}{a}$ et $\gamma$ fierent quantitates negativae. Quia vero $A$ et $A'$ signa eadem habent, $\sqrt{D}$ cadet inter $\sqrt{D + \frac{AA'}{aa'}}$ et $\sqrt{D - \frac{AA'}{aa'}}$ adeoque in hocce casu $-\frac{\gamma}{\alpha}$ inter $-\frac{\sqrt{D} + b}{a}$ et $-\frac{\sqrt{D} - b}{a}$. Quare pars prima theorematis pro casu priori est demonstrata.

Eodem modo perspicitur, in (5) et (6) necessario signa inferiora accipi debere quando neque $-\frac{\sqrt{D} + b}{a'}$ neque $-\frac{\sqrt{D} - b}{a'}$ signum idem habeant ut $a'$ sive $a$ quia accepto superiori $-\frac{\delta}{\beta}$ necessario fierent quantitates positivae. Unde protinus sequitur, $-\frac{\delta}{\beta}$ pro hocce casu iacere inter $-\frac{\sqrt{D} + b}{a'}$ et $-\frac{\sqrt{D} - b}{a'}$. Demonstrata est itaque etiam pars secunda theorematis pro casu posteriori. Quodsi aeque facile ostendi posset, in (3) et (4) signa inferiora accipi debere, quando neutra quantitatum $-\frac{\gamma}{\alpha}$, $-\frac{\delta}{\beta}$ signum idem habeat ut $a$, et in (5) et (6) superiora, quando neque $-\frac{\delta}{\beta}$ neque $-\frac{\gamma}{\alpha}$ signum oppositum habeat: hinc simili modo sequeretur, pro illo casu $-\frac{\gamma}{\alpha}$ iacere inter $-\frac{\sqrt{D} + b}{a}$ et $-\frac{\sqrt{D} - b}{a}$, pro hoc $-\frac{\delta}{\beta}$ inter $-\frac{\sqrt{D} + b}{a'}$ et $-\frac{\sqrt{D} - b}{a'}$ sive pars prima theorematis etiam pro casu posteriori et secunda pro casu priori demonstratae forent. Sed quum illud difficile quidem non sit, attamen sine quibusdam ambagibus fieri nequeat, methodum sequentem praeferimus.

Quando nullus numerorum $\alpha$, $\beta$, $\gamma$, $\delta = 0$, $-\frac{\gamma}{\alpha}$ et $-\frac{\delta}{\beta}$ eadem signa habebunt ut $-\frac{\sqrt{D} + b}{a}$ vel $-\frac{\sqrt{D} - b}{a}$. Quando itaque neutra harum quantitatum signum idem habet ut $a'$ sive $a$, adeoque $-\frac{\gamma}{\alpha}$ et $-\frac{\delta}{\beta}$ inter $-\frac{\sqrt{D} + b}{a}$ et $-\frac{\sqrt{D} - b}{a}$ cadit: neutra quantitatum $-\frac{\gamma}{\alpha}$, $-\frac{\delta}{\beta}$ signum idem ut $a$ habebit, cadetque $-\frac{\delta}{\beta} = \frac{-a'\gamma}{a\beta}$ (propter $aa' = DD - bb$) inter $-\frac{\sqrt{D} + b}{a}$ et $-\frac{\sqrt{D} - b}{a}$. Quare pro eo casu, ubi neque $\alpha$ neque $\beta = 0$, pars prima theor.~etiam pro casu secundo est demonstrata (nam conditio, ut neque $\gamma$ neque $\delta = 0$, iam in theor.~ipso est adiecta). Simili modo, quando nullus numerorum $\alpha$, $\beta$, $\gamma$, $\delta = 0$, et neque $-\frac{\gamma}{\alpha}$ neque $-\frac{\delta}{\beta}$ signum signo ipsius $a$ vel $a'$ oppositum habet adeoque $-\frac{\gamma}{\alpha}$ inter $-\frac{\sqrt{D} + b}{a}$ et $-\frac{\sqrt{D} - b}{a}$ iacet: etiam $\frac{\delta}{\beta}$ et $\frac{\gamma}{\alpha}$ signum oppositum signo ipsius $a'$ non habebit, cadetque $-\frac{\delta}{\beta} = \frac{-a\gamma}{a'\beta}$ inter $-\frac{\sqrt{D} + b}{a'}$ et $-\frac{\sqrt{D} - b}{a'}$. In eo igitur casu, ubi neque $\gamma$ neque $\delta = 0$, pars secunda theor.~etiam pro casu secundo est demonstrata.

Nihil itaque superesset quam ut demonstretur, partem primam theor.~etiam pro casu secundo locum habere, si alteruter numerorum $\alpha$, $\beta$ sit $= 0$, et partem secundam pro casu primo si aut $\gamma$ aut $\delta = 0$. At omnes hi casus sunt impossibiles.

Supponamus enim, pro parte prima theor., esse neque $\gamma$ neque $\delta = 0$; $-\frac{\gamma}{\alpha}$ non habere signum idem ut $a$ atque esse 1) $\alpha = 0$. Tum ex aequ.~$\alpha\delta - \beta\gamma = +1$ fit $\beta = +1$, $\gamma = +1$. Hinc ex (1) $A = -a'$, quare $A > a'$, adeoque etiam $a$ et $A$ signa opposita habent, unde fit $\sqrt{D - \frac{AA'}{aa'}} > \sqrt{D} > b$. Hinc patet, in (4) necessario signum inferius accipi debere, quia accepto superiori $-\frac{\delta}{\beta}$ manifesto signum idem obtineret ut $a'$. Fit itaque $-\frac{\delta}{\beta} > \frac{-\sqrt{D} + b}{-a'} > \frac{\sqrt{D}}{a'}$ (propter $a' < \sqrt{D} + b$ ex def.~formae reductae), Q.E.A.~quia $\frac{\delta}{\beta} = +1$ et $\delta$ non $= 0$. Sit $\beta = +1$, 2) $\alpha = 0$. Tum ex aequ.~$\alpha\delta - \beta\gamma = +1$ fit $\alpha = +1$, $\delta = +1$. Hinc ex (2) $-A' = a'$, quare $a'$ et $a$ et $A$ signa eadem habebunt, unde fit $\sqrt{D + \frac{AA'}{aa'}} > \sqrt{D} > b$. Hinc patet, in (3) signum inferius accipi debere, quia accepto superiori $-\frac{\gamma}{\alpha}$ signum idem obtineret ut $a$. Fit itaque $-\frac{\gamma}{\alpha} > \frac{-\sqrt{D} + b}{a} > 1$, Q.E.A.~eadem ratione ut ante.

Pro parte secunda si supponimus neque $\alpha$ neque $\beta = 0$; $-\frac{\delta}{\beta}$ non habere signum signo ipsius $a'$ oppositum atque 1) $\gamma = 0$: ex aequ.~$\alpha\delta - \beta\gamma = +1$ fit $\alpha = +1$, $\delta = +1$. Hinc ex (1) $A = a$, quare $a'$ et $A$ signa eadem habebunt, unde fit $\sqrt{D + \frac{AA'}{aa'}} > \sqrt{D} > b$. Quocirca in (6) signum superius erit accipiendum, quia accepto inferiori, $\frac{\delta}{\beta}$ obtineret signum oppositum signo ipsius $a'$. Fit igitur $\frac{\delta}{\beta} < \frac{-\sqrt{D} + b}{-a'} < 1$, Q.E.A., quia $\delta = +1$ et $\beta$ non $= 0$. Tandem 2) si esset $\delta = 0$, ex $\alpha\delta - \beta\gamma = +1$ fit $\beta = +1$, $\gamma = +1$, adeoque ex (2) $-A' = a'$. Hinc $\sqrt{D - \frac{AA'}{aa'}} > \sqrt{D} > b$, quare in (5) signum superius accipiendum. Hinc $\frac{\gamma}{\alpha} < \frac{-\sqrt{D} + b}{a'} < 1$. Quare theorema in omni sua extensione est demonstratum.

Quum differentia inter $-\frac{\sqrt{D} + b}{a}$ et $-\frac{\sqrt{D} - b}{a}$ sit $= \frac{2b}{a} < 1$: differentia inter $-\frac{\gamma}{\alpha}$ et $-\frac{\sqrt{D} + b}{a}$ vel $-\frac{\sqrt{D} - b}{a}$ erit $< \frac{1}{|\alpha|}$; inter $-\frac{\delta}{\beta}$ autem et $-\frac{\sqrt{D} + b}{a'}$, vel inter illam quantitatem et $-\frac{\sqrt{D} - b}{a'}$ nulla fractio iacere poterit, cuius denominator non sit maior quam $\gamma$ aut $\delta$ (lemma prae.). Eodem modo differentia quantitatis $-\frac{\delta}{\beta} = \frac{-a\gamma}{a'\beta}$ a fractione $-\frac{\gamma}{\alpha}$ vel hac a $\frac{\delta}{\beta}$ erit minor quam $\frac{1}{|\beta|}$, et inter illam quantitatem et $-\frac{\gamma}{\alpha}$ neutram harum fractionum iacere potest fractio, cuius denominator non sit maior quam $\alpha$ et $\beta$.

\bigskip

\textbf{192.}

Ex applicatione theor.~prae.~ad algorithmum art.~188 sequitur, quantitatem $-\frac{\gamma}{\alpha}$, quam per $L$ designabimus, iacere inter $-\frac{\gamma'}{\alpha'}$ et $-\frac{\gamma''}{\alpha''}$; inter $-\frac{\gamma''}{\alpha''}$ et $-\frac{\gamma'''}{\alpha'''}$; inter $-\frac{\gamma'''}{\alpha'''}$ et $-\frac{\gamma^{IV}}{\alpha^{IV}}$ etc.~(facile enim ex art.~189,~3 fin.~deducitur, nullum horum limitum habere signum oppositum signo ipsius $a$; quare quantitati radicali $\sqrt{D}$ signum positivum tribui debet) sive inter $-\frac{\gamma'}{\alpha'}$ et $-\frac{\gamma''}{\alpha''}$; inter $-\frac{\gamma''}{\alpha''}$ et $-\frac{\gamma'''}{\alpha'''}$ etc.

Omnes itaque fractiones $-\frac{\gamma'}{\alpha'}$, $-\frac{\gamma'''}{\alpha'''}$, $-\frac{\gamma^V}{\alpha^V}$ etc.~ab eadem parte iacebunt, omnesque $-\frac{\gamma''}{\alpha''}$, $-\frac{\gamma^{IV}}{\alpha^{IV}}$, $-\frac{\gamma^{VI}}{\alpha^{VI}}$ etc.~ipsi $-\frac{\gamma}{\alpha}$ a parte altera. Quoniam vero $\gamma' < \gamma'' < \gamma'''$ etc., $-\frac{\gamma''}{\alpha''}$ iacebit extra $-\frac{\gamma'}{\alpha'}$ et $-\frac{\gamma'''}{\alpha'''}$, similique ratione $-\frac{\gamma^{IV}}{\alpha^{IV}}$ extra $-\frac{\gamma'''}{\alpha'''}$ et $-\frac{\gamma^V}{\alpha^V}$ etc.~Unde manifestum est, has quantitates iacere sequenti ordine
\[
-\frac{\gamma'}{\alpha'} > -\frac{\gamma'''}{\alpha'''} > -\frac{\gamma^V}{\alpha^V} \; \ldots \; > L \; > \; \ldots \; > -\frac{\gamma^{VI}}{\alpha^{VI}} > -\frac{\gamma^{IV}}{\alpha^{IV}} > -\frac{\gamma''}{\alpha''}.
\]
Differentia autem inter $-\frac{\gamma}{\alpha}$ et $L$ erit minor quam differentia inter $-\frac{\gamma'}{\alpha'}$ et $-\frac{\gamma''}{\alpha''}$, i.e.~$< \frac{1}{|\alpha'\alpha''|}$; similique ratione differentia inter $-\frac{\gamma''}{\alpha''}$ et $L$ erit $< \frac{1}{|\alpha''\alpha'''|}$ etc.~Quamobrem $-\frac{\gamma'}{\alpha'}$, $-\frac{\gamma''}{\alpha''}$, $-\frac{\gamma'''}{\alpha'''}$ etc.~continuo propius ad limitem $L$ accedunt, et quoniam $\alpha'$, $\alpha''$, $\alpha'''$ continuo in infinitum crescunt, differentia fractionum a limite quavis quantitate data minor fieri potest.

Ex art.~189 nulla quantitatum $'\gamma/'\alpha$, $''\gamma/''\alpha$, $'''\gamma/'''\alpha$ etc.~signum idem habebit ut $a$: hinc per ratiocinia praecedentibus omnino similia sequitur, illas et hanc $'\gamma/'\alpha$, quam per $L'$ designabimus, iacere sequenti ordine
\[
\frac{'\alpha}{'\gamma} > \frac{''\alpha}{''\gamma} > \frac{'''\alpha}{'''\gamma} \; \ldots \; > L' \; > \; \ldots \; > \frac{'''\alpha}{'''\gamma} > \frac{''\alpha}{''\gamma} > \frac{'\alpha}{'\gamma}.
\]
Differentia autem inter $-\frac{'\gamma}{'\alpha}$ et $L'$ minor erit quam $\frac{1}{|'\alpha \; ''\alpha|}$, differentia inter $-\frac{''\gamma}{''\alpha}$ et $L'$ minor quam $\frac{1}{|''\alpha \; '''\alpha|}$ etc.~Quare fractiones $'\gamma/'\alpha$, $''\gamma/''\alpha$, $'''\gamma/'''\alpha$ etc.~continuo propius ad $L'$ accedent, et differentia quavis quantitate data minor fieri poterit.

\textit{In ex.~art.~188} fit $L = -\frac{\gamma'''}{\alpha'''} = -\frac{143}{805} = 0{,}1776388$ et fractiones appropinquantes $-\frac{\gamma}{\alpha} = 0$, $-\frac{\gamma'}{\alpha'} = \frac{1}{5}$, $-\frac{\gamma''}{\alpha''} = \frac{2}{11}$, $-\frac{\gamma'''}{\alpha'''} = \frac{3}{16}$, $-\frac{\gamma^{IV}}{\alpha^{IV}} = \frac{5}{27}$ etc.~Est autem $L = 0{,}1776388$\ldots

Ibidem fit $L' = -\frac{'\gamma}{'\alpha} = \frac{27}{152} = 0{,}1776315$, fractionesque approximantes $\frac{'\gamma}{'\alpha} = -\frac{1}{5}$, $-\frac{''\gamma}{''\alpha} = -\frac{2}{11}$, $-\frac{'''\gamma}{'''\alpha} = -\frac{3}{16}$, $-\frac{''''\gamma}{''''\alpha} = -\frac{5}{27}$, $-\frac{'''''\gamma}{'''''\alpha} = -\frac{8}{43}$ etc.~Est vero $L' = 0{,}1776315$\ldots

\bigskip

\textbf{193.}

\textbf{Theorem.} Si formae reductae $f$, $F$ proprie aequivalentes sunt: altera in alterius periodo contenta erit.

Sit $f = (a, b, -a')$, $F = (A, B, -A')$, determinans harum formarum $D$, transeatque illa in hanc per substitutionem propriam $\mathfrak{A}$, $\mathfrak{B}$, $\mathfrak{C}$, $\mathfrak{D}$. Tum dico, si periodus formae $f$ quaeratur progressioque utrimque infinita formarum reductarum atque transformationum formae $f$ in ipsas eruatur, eodem modo ut art.~188: vel $+\mathfrak{A}$ fore aequalem termino alicui progressionis $\ldots''\alpha$, $'\alpha$, $\alpha$, $\alpha'$, $\alpha''$\ldots, hocque posito $= \alpha^m$, $+\mathfrak{B}$ fore $= \beta^m$, $+\mathfrak{C} = \gamma^m$, $+\mathfrak{D} = \delta^m$; vel $-\mathfrak{A}$ fore aequalem termino alicui $'\alpha^m$, $-\mathfrak{B} = '\beta^m$, $-\mathfrak{C} = '\gamma^m$, $-\mathfrak{D} = '\delta^m$ (ubi $m$ etiam indicem negativum designare potest). In utroque casu $F$ manifesto identica erit cum $f^m$.

\textit{Dem.} I. Habentur quatuor aequationes
\begin{align*}
a\mathfrak{A}\mathfrak{A} + 2b\mathfrak{A}\mathfrak{C} - a'\mathfrak{C}\mathfrak{C} &= A \tag{1}\\
a\mathfrak{A}\mathfrak{B} + b(\mathfrak{A}\mathfrak{D} + \mathfrak{B}\mathfrak{C}) - a'\mathfrak{C}\mathfrak{D} &= B \tag{2}\\
a\mathfrak{B}\mathfrak{B} + 2b\mathfrak{B}\mathfrak{D} - a'\mathfrak{D}\mathfrak{D} &= -A' \tag{3}\\
\mathfrak{A}\mathfrak{D} - \mathfrak{B}\mathfrak{C} &= 1 \tag{4}
\end{align*}

Consideramus autem primo casum, ubi aliquis numerorum $\mathfrak{A}$, $\mathfrak{B}$, $\mathfrak{C}$, $\mathfrak{D} = 0$.

$1^\circ$ Si $\mathfrak{D} = 0$, fit ex (4) $\mathfrak{B}\mathfrak{C} = -1$, adeoque $\mathfrak{B} = +1$, $\mathfrak{C} = +1$. Hinc ex (1), $-a' = A$; ex (2), $-b - a' = B$ sive $B \equiv -b \pmod{A}$; unde sequitur formam $(A, B, -A')$ formae $(a, b, -a')$ ab ultima parte contiguam esse. Quoniam vero illa est reducta, necessario cum $f'$ identica erit. Ergo $B = b'$, adeoque ex (2) $b + b' = -a'\mathfrak{B}\mathfrak{D} = +a'\mathfrak{D}$; hinc propter $\mathfrak{A}\mathfrak{D} = A'$ fit $\mathfrak{D} = h'\mathfrak{A}'$. Unde colligitur, $+\mathfrak{A}$, $+\mathfrak{B}$, $+\mathfrak{C}$, $+\mathfrak{D}$ esse resp.~$= 0$, $-1$, $+1$, $h'$ sive, $= \alpha'$, $\beta'$, $\gamma'$, $\delta'$.

$2^\circ$ Si $\mathfrak{B} = 0$, fit ex (4) $\mathfrak{A} = \pm 1$, $\mathfrak{D} = \pm 1$; ex (3) $a' = A'$; ex (2) $b \pm a'\mathfrak{A} = B$, sive $b \equiv B \pmod{a'}$. Quoniam vero tum $b$ tum $B$ iacebunt inter $\sqrt{D}$ et $\sqrt{D} \mp a'$ (prout $a'$ pos.~vel neg., art.~184,~5). Quare erit necessario $b = B$, et $\mathfrak{D} = 0$. Hinc formae $f$, $F$ sunt identicae atque $+\mathfrak{A}$, $+\mathfrak{B}$, $+\mathfrak{C}$, $+\mathfrak{D} = 1$, $0$, $0$, $1 = \alpha$, $\beta$, $\gamma$, $\delta$ (resp.).

$3^\circ$ Si $\mathfrak{C} = 0$, fit ex (4) $\mathfrak{A} = +1$, $\mathfrak{D} = +1$; ex (1) $a = A$; ex (2) $+ a\mathfrak{B} + b = B$ sive $b \equiv B \pmod{a}$. Quia vero tum $b$ tum $B$ iacent inter $\sqrt{D}$ et $\sqrt{D} \mp a$: erit necessario $B = b$ et $\mathfrak{B} = 0$. Quare casus hic a praecedente non differt.

$4^\circ$ Si $\mathfrak{A} = 0$, ex (4) $\mathfrak{B} = +1$, $\mathfrak{C} = +1$; ex (3) $a' = -A''$, ex (2) $+ b\mathfrak{B} - a' = B$ sive $B \equiv -b \pmod{a'}$. Hinc forma $F$ formae $f$ a parte prima contigua erit, et proin cum forma $'f$ identica. Quare propter $'f = F$, et $B = 'b$, erit $+\mathfrak{D} = A$. Unde colligitur $+\mathfrak{A}$, $+\mathfrak{B}$, $+\mathfrak{C}$, $+\mathfrak{D}$ resp.~esse $= h$, $1$, $-1$, $0$, $= '\alpha$, $\beta'$, $\gamma'$, $\delta'$.

Superest itaque casus, ubi nullus numerorum $\mathfrak{A}$, $\mathfrak{B}$, $\mathfrak{C}$, $\mathfrak{D} = 0$. Hic per Lemma art.~190 quantitates $\mathfrak{A}$, $\mathfrak{B}$, $\mathfrak{C}$, $\mathfrak{D}$ idem signum habebunt, oriunturque inde duo casus, quum signum hoc vel cum signo ipsorum $a$, $a'$ convenire vel ipsi oppositum esse possit.

II. Si $\mathfrak{C}$, $\mathfrak{D}$ idem signum habent ut $a$: quantitas $-\frac{\mathfrak{C}}{\mathfrak{A}}$ (quam designabimus per $L$) inter has fractiones $-\frac{\gamma'}{\alpha'}$, $-\frac{\gamma''}{\alpha''}$, $-\frac{\gamma'''}{\alpha'''}$ etc.~sita erit (art.~191). Demonstrabimus iam, $\frac{\mathfrak{C}}{\mathfrak{A}}$ aequalem fore alicui fractionum $-\frac{\gamma'}{\alpha'}$, $-\frac{\gamma''}{\alpha''}$, $-\frac{\gamma'''}{\alpha'''}$ etc., atque $\frac{\mathfrak{D}}{\mathfrak{B}}$ proxime sequenti, scilicet si $-\frac{\mathfrak{C}}{\mathfrak{A}}$ fuerit $= -\frac{\gamma^m}{\alpha^m}$, $\frac{\mathfrak{D}}{\mathfrak{B}}$ fore $= -\frac{\delta^m}{\beta^m}$. In art.~prae.~ostendimus, quantitates $-\frac{\gamma'}{\alpha'}$, $-\frac{\gamma''}{\alpha''}$, $-\frac{\gamma'''}{\alpha'''}$ etc.~(quas brevitatis gratia per (1), (2), (3) etc.~denotabimus), atque $L$, hunc ordinem (I) observare: $(1)$, $(3)$, $(5)$\ldots$L$\ldots$(6)$, $(4)$, $(2)$; prima harum quantitatum est $= 0$ (propter $\alpha' = 0$), reliquae omnes idem signum habent ut $L$ sive $a$. Quoniam vero per hyp.~$\mathfrak{C}$, $\mathfrak{D}$ (pro quibus scribemus $\mathfrak{M}$, $\mathfrak{N}$) idem signum habent: patet has quantitates ipsi (1) a dextra iacere (aut si mavis ab eadem parte a qua $L$), et quidem, quum $L$ iaceat inter ipsas, alteram ipsi $L$ a dextra, alteram a laeva. Facile vero ostendi potest, $\mathfrak{M}$ ipsi (2) a dextra iacere non posse, alioquin enim $\mathfrak{M}$ iaceret inter (1) et $L$, unde sequeretur primo (2) iacere inter $\mathfrak{M}$ et $\mathfrak{N}$, adeoque denominatorem fractionis (2) maiorem esse denominatore fractionis $\mathfrak{M}$ (art.~190), secundo $\mathfrak{N}$ iacere inter (1) et (2), adeoque denom.~fractionis $\mathfrak{N}$ esse maiorem quam denom.~fractionis (2). Q.E.A.

Supponamus $\mathfrak{M}$ nulli fractionum (2), (3), (4) etc.~aequalem esse, ut, quid inde sequatur, videamus. Tum manifestum est, si fractio $\mathfrak{M}$ ipsi $L$ a laeva iaceat, necessario eam sitam esse aut inter (1) et (3), aut inter (3) et (5), aut inter (5) et (7) etc.~(quoniam $L$ est irrationalis, adeoque ipsi $\mathfrak{M}$ certo inaequalis, fractionesque (1), (3), (5) etc.~quavis quantitate data, ipsi $L$ inaequali propius ad $L$ accedere possunt). Si vero $\mathfrak{M}$ ipsi $L$ a dextra iacet: necessario iacebit aut inter (2) et (4), aut inter (4) et (6) aut inter (6) et (8) etc. Ponamus itaque $\mathfrak{M}$ iacere inter $(m)$ et $(m+2)$, patetque quantitates $\mathfrak{M}$, $(m)$, $(m+1)$, $(m+2)$, $L$ iacere sequenti ordine
\[
\text{(II)\footnote{Nihil hic refert, sive ordo in (II) idem sit ut in (I), sive huic oppositus, i.e.~sive $(m)$ etiam in (I) ipsi $L$ a laeva iaceat sive a dextra.}}: \quad \mathfrak{M}, \; (m), \; (m+2), \; L, \; (m+1).
\]
Tum erit necessario $\mathfrak{N} = (m+1)$. Iacebit enim $\mathfrak{N}$ ipsi $L$ a dextra; si vero etiam ipsi $(m+1)$ a dextra iaceret, $(m+1)$ iaceret inter $\mathfrak{N}$ et $\mathfrak{M}$, unde $\gamma^{m+1} > \mathfrak{N}$, $\mathfrak{M}$ vero inter $(m)$ et $(m+1)$, unde $\gamma^{m+1} > \mathfrak{M}$ (art.~190), Q.E.A.; si vero $\mathfrak{N}$ ipsi $(m+1)$ a laeva iaceret, sive inter $(m+2)$ et $(m+1)$, foret $\mathfrak{D} > \gamma^{m+1}$ et quia $(m+2)$ inter $\mathfrak{M}$ et $\mathfrak{N}$ foret $\gamma^{m+1} > \mathfrak{M}$, Q.E.A. Erit itaque $\mathfrak{N} = (m+1)$ sive $\frac{\mathfrak{D}}{\mathfrak{B}} = -\frac{\delta^m}{\beta^m} = \frac{\gamma^{m+1}}{\alpha^{m+1}}$.

Quia $\mathfrak{A}\mathfrak{D} - \mathfrak{B}\mathfrak{C} = 1$, $\mathfrak{A}$ erit primus ad $\mathfrak{C}$ et ex simili ratione $\mathfrak{B}$ primus ad $\mathfrak{D}$. Unde facile perspicitur aequationem $\frac{\mathfrak{C}}{\mathfrak{A}} = \frac{\gamma^m}{\alpha^m}$ consistere non posse, nisi fuerit aut $\mathfrak{A} = \alpha^m$, $\mathfrak{D} = \delta^m$, aut $\mathfrak{A} = -\alpha^m$, $\mathfrak{D} = -\delta^m$. Iam quum forma $f$ per substitutionem propriam $\alpha^m$, $\beta^m$, $-\gamma^m$, $-\delta^m$ in formam $f^m$ transmutetur, quae est $(+a^m, b^m, +a^{m+1})$: habebuntur aequationes
\begin{align*}
a\alpha^m\alpha^m + 2b\alpha^m\gamma^m - a'\gamma^m\gamma^m &= +a^{m+1} \tag{5}\\
a\alpha^m\beta^m + b(\alpha^m\delta^m + \beta^m\gamma^m) - a'\gamma^m\delta^m &= b^m \tag{6}\\
a\beta^m\beta^m + 2b\beta^m\delta^m - a'\delta^m\delta^m &= +a^{m+2} \tag{7}\\
\alpha^m\delta^m - \beta^m\gamma^m &= 1 \tag{8}
\end{align*}

Hinc fit: (ex aequ.~7 et 3), $\pm a^{m+2} = A$. Porro multiplicando aequationem (2) per $a\alpha^m\delta^m - \beta^m\gamma^m$, aequationem (6) per $2(\mathfrak{A}\mathfrak{D} - \mathfrak{B}\mathfrak{C})$ et subtrahendo facile per evolutionem confirmatur esse
\[
B - b^m = (\mathfrak{C}\alpha^m - \mathfrak{A}\gamma^m)(a\mathfrak{B}\beta^m + b(\mathfrak{A}\delta^m + \mathfrak{D}\alpha^m) - \beta'\mathfrak{D}\gamma^m) + b((\mathfrak{C}\alpha^m + \mathfrak{A}\gamma^m) - \alpha(\mathfrak{C}\gamma^m)) + (\mathfrak{B}\delta^m - \mathfrak{D}\beta^m)(a\mathfrak{A}\alpha^m) \tag{9}
\]
sive quoniam vel $\mathfrak{A}\alpha^m = \mathfrak{D}\delta^m$, $\mathfrak{C}\gamma^m = \mathfrak{D}\beta^m$ vel $\mathfrak{A}\alpha^m = -\mathfrak{D}\delta^m$, $\mathfrak{C}\gamma^m = -\mathfrak{D}\beta^m$,
\[
B - b^m = \pm((\mathfrak{C}\alpha^m - \mathfrak{A}\gamma^m)(\mathfrak{B}\delta^m - \mathfrak{D}\beta^m) \cdot 2b\gamma^m) \pm a'\mathfrak{D}\gamma^m \cdot 2\gamma^m.
\]
Hinc $B \equiv b^m \pmod{A'}$; quia vero tum $B$ tum $b^m$, inter $\sqrt{D}$ et $\sqrt{D} \mp A$ iacent, necessario erit $B = b^m$, adeoque $\mathfrak{C}\alpha^m - \mathfrak{A}\gamma^m = 0$, sive $\frac{\mathfrak{C}}{\mathfrak{A}} = \frac{\gamma^m}{\alpha^m}$, i.e.~$\mathfrak{M} = (m)$.

Hoc modo itaque ex suppositione, $\mathfrak{M}$ nulli quantitatum (2), (3), (4) etc.~aequalem esse, deduximus, eam revera alicui aequalem esse. Quodsi vero ab initio supponimus, esse $\frac{\mathfrak{C}}{\mathfrak{A}} = \frac{\gamma^m}{\alpha^m}$ vel $= (m)$, manifesto erit vel $\mathfrak{A} = \alpha^m$, $\mathfrak{C} = \gamma^m$, $-\mathfrak{A} = \alpha^m$, $-\mathfrak{C} = \gamma^m$. In utroque casu fit ex (1) et (5) $A = \pm a^m$, et ex (9) $B - b^m = \pm(\mathfrak{B}\delta^m - \mathfrak{D}\beta^m) \cdot 2b\gamma^m$ sive $B \equiv b^m \pmod{A}$. Hinc simili modo ut supra concluditur $B = b^m$, et hinc $\mathfrak{B}\delta^m - \mathfrak{D}\beta^m = 0$; quare quum $\mathfrak{B}$ ad $\mathfrak{D}$ primus sit et $\beta^m$ ad $\delta^m$: erit aut $\mathfrak{B} = \beta^m$, $\mathfrak{D} = \delta^m$, aut $-\mathfrak{B} = \beta^m$, $-\mathfrak{D} = \delta^m$, et proin ex (7) $-A' = \pm a^{m+2}$. Quamobrem formae $F$, $f^m$ identicae erunt. Adiumento aequationis $\mathfrak{A}\mathfrak{D} - \mathfrak{B}\mathfrak{C} = \alpha^m\delta^m - \beta^m\gamma^m = 1$ autem nullo negotio probatur, poni debere $-\frac{\mathfrak{B}}{\mathfrak{A}} = \beta^m$, $-\frac{\mathfrak{D}}{\mathfrak{C}} = \delta^m$, quando $+\mathfrak{A} = \alpha^m$, $+\mathfrak{C} = \gamma^m$; contra $-\mathfrak{B} = \beta^m$, $-\mathfrak{D} = \delta^m$, quando $-\mathfrak{A} = \alpha^m$, $-\mathfrak{C} = \gamma^m$. \hfill Q.E.D.

III. Si signum quantitatum $\mathfrak{C}$, $\mathfrak{D}$ signo ipsius $a$ oppositum: demonstratio praecedenti tam similis est, ut praecipua tantum momenta addigavisse sufficiat. $-\frac{\mathfrak{C}}{\mathfrak{A}}$ iacebit inter $\frac{'\gamma}{'\alpha}$ et $\frac{''\gamma}{''\alpha}$. Fractio $\frac{\mathfrak{C}}{\mathfrak{A}}$ alicui fractionum $'\gamma/'\alpha$, $''\gamma/''\alpha$, $'''\gamma/'''\alpha$ etc.~aequalis erit (I); qua posita $-\frac{\mathfrak{D}}{\mathfrak{B}} = '\delta^m/'\beta^m$, $-\frac{\mathfrak{D}}{\mathfrak{C}}$ erit $= '\delta^m/'\gamma^m$ (II).

Demonstratur autem (I) ita: Si $\frac{\mathfrak{C}}{\mathfrak{A}}$ nulli illarum fractionum aequalis esse supponitur: inter duas tales $\frac{'\gamma}{'\alpha}$ et $\frac{^{m+1}\gamma}{^{m+1}\alpha}$ iacere debebit. Hinc vero eodem modo ut supra deducitur, necessario esse $-\frac{\mathfrak{D}}{\mathfrak{B}} = \frac{'\delta}{'\beta}$ atque vel $\mathfrak{A} = '\alpha$, $\mathfrak{C} = '\gamma$, vel $-\mathfrak{A} = '\alpha$, $-\mathfrak{C} = '\gamma$. Quoniam vero $f$ per substitutionem propriam $'\alpha$, $'\beta$, $'\gamma$, $'\delta$ in formam $'f$ transit: hinc emergunt tres aequationes, ex quibus coniunctis cum aequ.~1, 2, 3, 4 atque hac, $'\alpha\cdot'\delta - '\beta\cdot'\gamma = 1$ deducitur eodem modo ut supra, terminum primum $A$ formae $F$ termino primo formae $'f$ aequalem esse, illiusque terminum medium medio huius congruum secundum modulum $A$, unde sequitur, quia utraque forma est reducta, adeoque utriusque terminus medius inter $\sqrt{D}$ et $\sqrt{D} \mp A$ situs, hos terminos medios aequales esse: hinc vero deducitur $\frac{\mathfrak{C}}{\mathfrak{A}} = '\gamma/'\alpha$. Veritas itaque assertionis (I) derivata hic est ex suppositione illam esse falsam.

Supponendo autem $\frac{\mathfrak{C}}{\mathfrak{A}} = '\gamma/'\alpha$, prorsus simili modo et per easdem aequationes demonstratur, $-\frac{\mathfrak{D}}{\mathfrak{B}} = '\delta/'\beta$ esse etiam $= '\delta/'\gamma$ quod erat secundum (II). Hinc vero adiumento aequationum $\mathfrak{A}\mathfrak{D} - \mathfrak{B}\mathfrak{C} = '\alpha\cdot'\delta - '\beta\cdot'\gamma = 1$ deducitur esse vel
\[
\mathfrak{A} = '\alpha, \; \mathfrak{B} = '\beta, \; \mathfrak{C} = '\gamma, \; \mathfrak{D} = '\delta
\]
vel
\[
\mathfrak{A} = -'\alpha, \; \mathfrak{B} = -'\beta, \; \mathfrak{C} = -'\gamma, \; \mathfrak{D} = -'\delta
\]
formasque $F$, $'f$ identicas. \hfill Q.E.D.

\bigskip

\textbf{194.}

Quum formae quas supra socias vocavimus (art.~187,~6), semper sint improprie aequivalentes (art.~159), perspicuum est, si formae reductae $F$, $f$ improprie aequivalentes sint, formaque $F$ socia forma $G$, formas $f$, $G$ proprie aequivalentes fore adeoque formam $G$ in periodo formae $f$ contentam. Quodsi itaque formae $F$, $f$ tum proprie tum improprie aequivalentes sunt, patet, tum $F$ tum $G$ in periodo formae $f$ reperiri debere. Quare periodus haec sibi ipsi socia erit, duasque formas ancipites continebit (art.~187,~7). Unde theorema art.~165 egregie confirmatur, ex quo iam poteramus esse certi, formam aliquam ancipitem dari formis $F$, $f$ aequivalentem.

\bigskip

\textbf{195.}

\textbf{Problema.} Propositis duabus formis quibuscunque $\Phi$, $\varphi$ eiusdem determinantis: diiudicare utrum aequivalentes sint, annon.

\textit{Sol.} Quaerantur duae formae reductae $F$, $f$ propositis $\Phi$, $\varphi$ resp.~proprie aequivalentes (art.~183). Quae prout aut proprie tantum aequivalent, aut improprie tantum, aut utroque modo, aut neutro; etiam propositae aut proprie tantum, aut improprie tantum, aut utroque aut neutro modo aequivalentes erunt. Evolvatur periodus alterutrius formae reductae e.g.~periodus formae $f$. Si forma $F$ in hac periodo occurrit neque vero simul forma ipsi $F$ socia, manifesto casus, primus locum habebit; contra si socia haec adest neque vero $F$ ipsa, secundus; si utraque, tertius; si neutra, quartus.

\textit{Ex.} Propositae sint formae $(129, 92, 65)$, $(42, 59, 81)$ determinantis $79$. His proprie aequivalentes inveniuntur reductae $(10, 7, -3)$, $(5, 8, -3)$. Periodus formae prioris haec est: $(10, 7, -3)$, $(-3, 8, 5)$, $(5, 7, -6)$, $(-6, 5, 9)$, $(9, 4, -7)$, $(-7, 3, 10)$. In qua quum forma $(5, 8, -3)$ ipsa non reperiatur, sed tamen socia $(-3, 8, 5)$: formas propositas improprie tantum aequivalere concludimus.

Si omnes formae reductae determinantis dati eodem modo ut supra (art.~187,~5) in periodos $P$, $Q$, $R$ etc.~distribuuntur, atque e quavis periodo forma aliqua ad libitum eligitur, ex $P$, $F$; ex $Q$, $G$; ex $R$, $H$ etc.: inter has formas $F$, $G$, $H$ etc.~duae quae proprie aequivalent esse non poterunt. Quaevis autem alia forma eiusdem determinantis alicui ex istis proprie aequivalens erit et quidem unicae tantum. Hinc manifestum est, omnes formas huius determinantis in totidem classes distribui posse, quot habeantur periodi, scilicet referendo eas, quae formae $F$ proprie aequivalent in primam classem, eas quae formae $G$ proprie aequivalent in secundam etc. Hoc modo omnes formae in eadem classe contentae proprie aequivalentes erunt, formae vero e classibus diversis non poterunt proprie aequivalere. Sed hic huic argumento infra fusius explicando non immoramur.

\bigskip

\textbf{196.}

\textbf{Problema.} Propositis duabus formis proprie aequivalentibus $\Phi$, $\varphi$: invenire transformationem propriam alterius in alteram.

\textit{Sol.} Per methodum art.~183 inveniri poterunt duae series formarum
\[
\Phi, \Phi', \Phi'' \ldots \Phi^n \quad \text{et} \quad \varphi, \varphi', \varphi'' \ldots \varphi^m
\]
tales ut quaevis forma sequens praecedenti proprie aequivaleat, ultimaeque $\Phi^n$, $\varphi^m$ sint formae reductae; et quum $\Phi$, $\varphi$ proprie aequivalentes esse supponantur, necessario $\Phi^n$ in periodo formae $\varphi^m$ contenta erit. Sit $\varphi^m = f$ ipsiusque periodus usque ad formam $\Phi^n$ haec
\[
f, \; f', \; f'' \; \ldots \; f^n = \Phi^n
\]
ita ut in hac periodo index formae $\Phi^n$ sit $m$; designenturque formae quae oppositae sunt sociis formarum
\[
\Phi, \; \Phi', \; \Phi'' \ldots \Phi^n \text{ per } \Psi, \; \Psi', \; \Psi'' \ldots \Psi^n \text{ resp.}\footnote{Ita ut $\Psi$ oriatur ex commutando terminum primum et ultimum tribuendoque medio signum oppositum, similiterque de reliquis.}
\]
Tum in progressione
\[
\varphi, \; \varphi', \; \varphi'' \; \ldots \; \varphi^m = f, \; f', \; f'' \; \ldots \; f^{n-1}, \; \Psi^{n-1}, \; \Psi^{n-2} \; \ldots \; \Psi', \; \Psi, \; \Phi
\]
quaevis forma praecedenti ab ultima parte contigua erit, unde per art.~177 inveniri poterit transformatio propria primae $\varphi$ in ultimam $\Phi$. Illud autem de formis reliquis progressionis nullo negotio perspicitur; de his $f^{n-1}$, $\Psi^{n-1}$ sic probatur: Sit
\[
f^{n-1} = (g, h, i); \quad f^n \text{ sive } \Phi^n = (g', h', i'); \quad f^{n+1} = (g'', h'', i'').
\]
Forma $(g', h', i')$ tum formae $(g, h, i)$ tum formae $(g'', h'', i'')$ ab ultima parte contigua erit; hinc $i = g' = i''$, et $-h \equiv h' \equiv -h'' \pmod{i \text{ vel } g \text{ vel } i''}$. Unde manifestum est, formam $(i'', -h'', g'')$, i.e.~formam $\Psi^{n-1}$ formae $(g, h, i)$, i.e.~formae $f^{n-1}$ ab ultima parte contiguam esse.

Si formae $\Phi$, $\varphi$ improprie aequivalentes sunt: forma $\varphi$ proprio aequivalebit formae, cui opposita est $\Phi$. Inveniri poterit itaque transformatio propria formae $\varphi$ in formam, cui est opposita; quae si supponitur fieri per substitutionem $\alpha$, $\beta$, $\gamma$, $\delta$, facile perspicitur, $\varphi$ improprie transformari in ipsam $\Phi$ per substitutionem $\alpha$, $-\beta$, $\gamma$, $-\delta$.

Hinc etiam perspicuum est, si formae $\Phi$, $\varphi$ tum proprie tum improprie aequivalentes sint, inveniri posse duas transformationes, propriam et impropriam.

\textit{Ex.} Quaeritur transformatio impropria formae $(129, 92, 65)$ in formam $(42, 59, 81)$, quam illi improprie aequivalere in art.~prae.~invenimus. Investiganda erit itaque primo transformatio propria formae $(129, 92, 65)$ in formam $(42, -59, 81)$. Ad hunc finem evolvitur progressio formarum haec: $(129, 92, 65)$, $(65, -27, 10)$, $(10, 7, -3)$, $(-3, 8, 5)$, $(5, 22, 81)$, $(81, 59, 42)$, $(42, -59, 81)$. Hinc deducitur transformatio propria $47$, $56$, $73$, $87$, per quam $(129, 92, 65)$ transit in $(42, -59, 81)$; quare per impropriam $-47$, $-56$, $73$, $87$ transibit in $(42, 59, 81)$.

\bigskip

\textbf{197.}

Si transformatio una formae alicuius $(a, h, c)\ldots \varphi$ in aequivalentem $\Phi$ habetur: ex hac omnes transformationes similes formae $\varphi$ in $\Phi$ deduci poterunt, si modo omnes solutiones aequationis indeterminatae $tt - Duu = mm$ assignari possunt, designante $D$ determinantem formarum $\Phi$, $\varphi$; $m$ divisorem communem maximum numerorum $a$, $2h$, $c$ (art.~162). Hoc igitur problema, quod pro valore negativo ipsius $D$ iam supra solvimus, nunc pro positivo aggrediemur. Quia vero manifesto quivis valor ipsius $t$ aequationi satisfaciens etiamnum mutato signo satisfacit, similiterque quivis valor ipsius $u$: sufficiet si omnes valores positivos ipsorum $t$, $u$ assignare possimus, fungeturque quaelibet solutio per valores positivos, quatuor solutionum vice. Hoc negotium ita absolvemus, ut primo valores minimos ipsorum $t$, $u$ (praeter hos per se obvios $t = m$, $u = 0$) invenire, tum ex his omnes reliquos derivare doceamus.

\bigskip

\textbf{198.}

\textbf{Problema.} Invenire numeros minimos $t$, $u$ aequationi indeterminatae $tt - Duu = mm$ satisfacientes, siquidem forma aliqua $(M, N, P)$ datur, cuius determinans est $D$, numerorumque $M$, $2N$, $P$ divisor communis maximus $= m$.

\textit{Sol.} Accipiatur ad libitum forma reducta $(a, b, -a')\ldots f$, determinantis $D$, ubi divisor communis maximus numerorum $a$, $2b$, $a'$ sit $m$, qualem dari vel inde manifestum est, quod forma reducta formae $(M, N, P)$ aequivalens inveniri potest, quae per art.~161 hac proprietate erit praedita; sed ad propositum praesens quaevis forma reducta, in qua conditio haec locum habet, poterit adhiberi.

Evolvatur periodus formae $f$, quam ex $n$ formis constare supponemus. Retentis omnibus signis, quibus in art.~188 usi sumus, $f^n$ erit $(+a^n, b^n, -a^{n+1})$, quia $n$ par, et in hanc formam transibit $f$ per substitutionem propriam $\alpha^n$, $\beta^n$, $\gamma^n$, $\delta^n$.

Quia vero $f$ et $f^n$ sunt identicae: $f$ transibit in $f^n$ etiam per substitutionem propriam $1$, $0$, $0$, $1$. Ex his duabus transformationibus similibus formae $f$ in $f^n$ per art.~162 deduci poterit solutio aequationis $tt - Duu = mm$ in integris, scilicet $t = \frac{1}{2}(\alpha^n + \delta^n)m$ (aequ.~18 art.~162), $u = \frac{\gamma^n}{m}$ (aequ.~19)\footnote{Quae in art.~162 erant $\alpha$, $\beta$, $\gamma$, $\delta$; $\alpha'$, $\beta'$, $\gamma'$, $\delta'$; $A$, $B$, $C$; $a$, $b$, $c$; $e$, $h$ etc.~hic sunt $1$, $0$, $0$, $1$; $\alpha^n$, $\beta^n$, $\gamma^n$, $\delta^n$; $a$, $b$, $-a'$; $a$, $b$, $-a'$; $1$.}. Designentur hi valores positive accepti si forte nondum sunt per $T$, $U$, eruntque hi $T$, $U$ valores minimi ipsorum $t$, $u$, praeter hos $t = m$, $u = 0$ (a quibus necessario erunt diversi, quia manifesto $\gamma^n$ non poterit esse $= 0$).

Supponamus enim dari adhuc minores valores ipsorum $t$, $u$ puta $t$, $u$ qui sint positivi et $u$ non $= 0$. Tum per art.~162 forma $f$ per substitutionem propriam $-(t - bu)$, $-\beta' u$, $-au$, $-(t + bu)$ transformabitur in formam cum ipsa identicam. Iam ex art.~193,~II sequitur, aut $-(t - bu)$ aut $+(t + bu)$ alicui numerorum $\alpha'$, $\alpha''$, $\alpha'''$ etc.~aequalem esse debere, puta $= \alpha^m$ (quia enim $u^2 = D\mu^2 + mm = bb\mu^2 + aa'\mu^2 + mm$, erit $u^2 > bb\mu^2$, adeoque $t > bu$ positivus; hinc fractio $-\frac{t - bu}{-au}$, quae respondet fractioni $-\frac{\gamma}{\alpha}$ in art.~193, idem signum habebit ut $a$ vel $a'$); atque in casu priori $-(t - bu)$, $-au$, $-(t + bu)$, in posteriori easdem quantitates mutatis signis, resp.~$= \alpha^m$, $\gamma^m$, $\beta^m$, $\delta^m$. Sed quum sit $u < U$ i.e.~$\mu < \frac{\gamma^n}{m}$ et $] > 0$: erit $\gamma^m < \gamma^n$ et $] > 0$; quocirca quum progressio $\gamma$, $\gamma'$, $\gamma''$ etc.~continuo crescat, necessario $\mu$ iacebit inter $0$ et $n$ excl. Forma vero respondens, $f^\mu$ identica erit cum forma $f$, Q.E.A., quum omnes formae $f$, $f'$, $f''$ etc.~usque ad $f^{n-1}$ diversae esse supponantur. Ex his colligitur, minimos valores ipsorum $t$, $u$ (exceptis valoribus $m$, $0$) esse $T$, $U$.

\textit{Ex.} Si $D = 79$, $m = 1$: adhiberi poterit forma $(3, 8, -5)$, pro qua $n = 6$, atque $\alpha^n = -805$, $\gamma^n = -143$, $\delta^n = -152$ (art.~188). Hinc $T = 80$, $U = 9$, qui sunt valores minimi numerorum $t$, $u$, aequationi $tt - 79uu = 1$ satisfacientes.

\bigskip

\textbf{199.}

Ad praxin formulae adhuc commodiores erui possunt. Erit nimirum $2b\gamma^n = -a(\alpha^n - \delta^n)$, quod facile ex art.~162 deducitur, multiplicando aequ.~[19] per $2b$, [20] per $a$ et mutando characteres illic adhibitos in praesentes. Hinc fit $\alpha^n + \delta^n = 2\delta^n - \frac{a'}{b}\gamma^n$, adeoque
\[
\pm T = \frac{m}{2}(\delta^n + \gamma^n), \quad \pm U = \frac{\gamma^n}{m}
\]
Per similem methodum bos valores obtinemus
\[
\pm T = \frac{m}{2}(\alpha^n + \delta^n), \quad \pm U = \frac{\alpha^n - \delta^n}{2b}
\]
Tum hae tum illae formulae perquam commode evadunt, propter $\gamma^n = \beta^{n-1}$, $\alpha^n = \delta^{n-1}$, ita ut si his uteris, solam progressionem $\delta'$, $\delta''$, $\delta'''$\ldots$\delta^n$; si illis uti mavis, solam hanc $\gamma'$, $\gamma''$, $\gamma'''$\ldots etc.~supputavisse sufficiat. Praeterea ex art.~189,~3 facile deducitur quum $n$ necessario sit par, $\alpha^n$ et $-\delta^n$ eadem signa habere; neque minus $\beta^n$ et $-\gamma^n$, ita ut in formula priori pro $T$ differentia absoluta, in posteriori summa absoluta accipi debeat neque adeo ad signa respicere omnino opus sit. Receptis signis in art.~189,~4 adhibitis erit ex formula priori
\[
T = m[k', k'', k'''\ldots k^{n-1}] - \frac{a'}{b}[k', k'', k'''\ldots k^{n-2}], \quad U = \frac{1}{m}[k', k'', k'''\ldots k^{n-1}]
\]
ex posteriori
\[
T = m[k', k''\ldots k^{n-1}] + \frac{a'}{b}[k'', k'''\ldots k^n], \quad U = \frac{1}{m}[k', k''\ldots k^n]
\]
umi pro valore ipsius $T$ etiam $m[k'', k'''\ldots k^n, -k']$ scribi poterit.

\textit{Ex.} Pro $D = 61$, $m = 1$ adhiberi potest forma $(2, 7, -6)$, pro qua eruitur $n = 6$; $k'$, $k''$, $k'''$, $k^{IV}$, $k^V$, $k^{VI}$ resp.~$= 2$, $2$, $7$, $2$, $2$, $7$. Hinc fit
\[
T = 2[2, 2, 7, 2, 2, 7] - 7[2, 2, 7, 2, 2] = 2888 - 1365 = 1523
\]
ex formula prima; idem provenit ex secunda
\[
T = 2[2, 7, 2, 2] + 7[2, 7, 2, 2, 7]
\]
$U$ vero fit $= [2, 2, 7, 2, 2] = \frac{1}{2}[2, 7, 2, 2, 7] = 195$.

Ceterum plura artificia adhuc dantur, per quae calculus contrahi potest, sed de his fusius hic loqui brevitas non permittit.

\bigskip

\textbf{200.}

Ut ex valoribus minimis ipsorum $t$, $u$ omnes obtineamus, aequationem $TT - DUU = mm$ ita exhibemus
\[
\left(\frac{T}{m} + \frac{U}{m}\sqrt{D}\right)\left(\frac{T}{m} - \frac{U}{m}\sqrt{D}\right) = 1
\]
inde etiam erit
\[
\left(\frac{T}{m} + \frac{U}{m}\sqrt{D}\right)^e \left(\frac{T}{m} - \frac{U}{m}\sqrt{D}\right)^e = 1 \tag{1}
\]
denotante $e$ numerum quemcunque. Iam designabimus brevitatis caussa valores quantitatum
\[
\frac{1}{2}m\left[\left(\frac{T}{m} + \frac{U}{m}\sqrt{D}\right)^e + \left(\frac{T}{m} - \frac{U}{m}\sqrt{D}\right)^e\right],
\]
\[
\frac{1}{2\sqrt{D}}m\left[\left(\frac{T}{m} + \frac{U}{m}\sqrt{D}\right)^e - \left(\frac{T}{m} - \frac{U}{m}\sqrt{D}\right)^e\right]
\]
generaliter per $t^e$, $u^e$ resp.~i.e.~illarum valores pro $e = 0$, per $t^0$, $u^0$ (qui erunt $m$, $0$); pro $e = 1$ per $t'$, $u'$ (qui fiunt $T$, $U$); pro $e = 2$ per $t''$, $u''$; pro $e = 3$ per $t'''$, $u'''$ etc.~demonstrabimusque, si pro $e$ accipiantur omnes numeri integri non negativi i.e.~$0$, omnesque positivi ab $1$ usque ad $\infty$, expressiones illas exhibere omnes valores positivos ipsorum $t$, $u$: scilicet I) omnes valores illarum expressionum esse revera valores ipsorum $t$, $u$; II) omnes valores illos esse numeros integros; III) nullos valores positivos ipsorum $t$, $u$ dari, qui sub formulis illis non contineantur\footnote{In his solis quatuor expressionibus et in aequ.~[1] $e$ denotat exponentem potestatis; in reliquis literae apici adscriptae semper indicem designant.}.

I. Substitutis pro $t^e$, $u^e$ valoribus suis nullo negotio adiumento aequ.~[1] confirmatur, esse
\[
(t^e + u^e\sqrt{D})(t^e - u^e\sqrt{D}) = mm, \quad \text{i.e.} \quad t^e t^e - Du^e u^e = mm.
\]

II. Eodem modo facile confirmatur, esse generaliter
\[
t^{e+1} = \frac{2T}{m}t^e - t^{e-1}, \quad u^{e+1} = \frac{2T}{m}u^e - u^{e-1}.
\]
Hinc manifestum est, duas progressiones $t^0$, $t'$, $t''$, $t'''$ etc., $u^0$, $u'$, $u''$, $u'''$ etc.~esse recurrentes, et utriusque scalam relationis $= \frac{2T}{m}$, $-1$, scilicet
\[
t'' = \frac{2T}{m}t' - t^0, \quad t''' = \frac{2T}{m}t'' - t' \text{ etc.}, \quad u'' = \frac{2T}{m}u' - u^0 \text{ etc.}
\]
Iam quoniam per hyp.~forma aliqua datur, $(M, N, P)$, determinantis $D$, in qua $M$, $2N$, $P$ per $m$ sunt divisibiles: habebitur
\[
TT = (NN - MP)UU + mm
\]
eritque adeo manifeste $\frac{2T}{m}$ per $m$ divisibilis. Hinc $\frac{2T}{m}$ erit numerus integer et quidem positivus. Quia vero $t^0 = m$, $t' = T$, $u^0 = 0$, $u' = U$, adeoque integri: omnes $t''$, $t'''$ etc., $u''$, $u'''$, etc.~etiam integri erunt. Porro perspicuum est, quia $TT > mm$, omnes $t^0$, $t'$, $t''$, $t'''$ etc.~positivos et continuo in infinitum crescentes esse, nec non omnes $u^0$, $u'$, $u''$, $u'''$ etc.

III. Supponamus, dari adhuc alios valores positivos ipsorum $t$, $u$, qui in progressione $t^0$, $t'$, $t''$ etc., $u^0$, $u'$, $u''$ etc.~non contenti sint, puta $\mathfrak{T}$, $\mathfrak{U}$. Manifestum est, quum progressio $u^0$, $u'$ etc.~a $0$ in infinitum crescat, $\mathfrak{U}$ necessario inter duos terminos proximos, $u^\lambda$ et $u^{\lambda+1}$ situm fore, ita ut sit $u^\lambda < \mathfrak{U} < u^{\lambda+1}$. Ut absurditatem huius suppositionis demonstremus, observamus

$1^\circ$ Aequationi $tt - Duu = mm$ satisfactum iri etiam ponendo
\[
t = \frac{1}{m}(\mathfrak{T}T \mp \mathfrak{U}UD), \quad u = \frac{1}{m}(\mathfrak{T}U \mp \mathfrak{U}T)
\]
unde, quoniam
\[
tt - Duu = \frac{1}{mm}(\mathfrak{T}\mathfrak{T} - D\mathfrak{U}\mathfrak{U})(TT - DUU) = \frac{1}{mm} \cdot mm \cdot mm = mm,
\]
hi valores ipsorum $t$, $u$ aequationi satisfacient.

$2^\circ$ Eosdem valores esse integros. Nam quum $TT - DUU = mm$, erit $T \equiv 0$, $U \equiv 0 \pmod{m}$, adeoque $\mathfrak{T}T \mp \mathfrak{U}UD$ et $\mathfrak{T}U \mp \mathfrak{U}T$ per $m$ divisibiles. Porro vero hi valores integri erunt, quoniam $t^\lambda t' - Du^\lambda u' = mm$ et $t^{\lambda+1}t' - Du^{\lambda+1}u' = mm$.

$3^\circ$ Valores istos esse minores quam $T$, $U$, contra hyp. Quoniam enim $u^\lambda < \mathfrak{U} < u^{\lambda+1}$ et $t^\lambda < \mathfrak{T}$, erit $\frac{1}{m}(\mathfrak{T}U - \mathfrak{U}T) < \frac{1}{m}(t^{\lambda+1}U - u^{\lambda+1}T) = U$; similique modo $\frac{1}{m}(\mathfrak{T}T - \mathfrak{U}UD) < T$. Q.E.A.

Quare omnes valores positivi ipsorum $t$, $u$ in progressione illa continentur, et problema penitus est solutum.

\end{document}
